% !TEX program = pdflatex
% Hugging Face Deep RL Course - Unit 2. This file reuses ../preamble.tex.
% Build from notes/en/: latexmk -pdf unit2en.tex (VS Code / LaTeX Workshop onSave).
\documentclass[11pt,a4paper]{article}

% Shared typography/layout for Hugging Face Deep RL Course notes.
% Compile with pdfLaTeX / latexmk -pdf with notes/ as the working directory.
% e.g. cd notes && latexmk -pdf unit1.tex

\usepackage[T1]{fontenc}
\usepackage[utf8]{inputenc}
\usepackage[english]{babel}
\usepackage{lmodern}
\usepackage{textcomp}
\usepackage{geometry}
\geometry{a4paper,margin=22mm,headheight=14pt,headsep=8mm,footskip=12mm}

\usepackage{amsmath,amssymb,mathtools}
\usepackage{graphicx,caption,needspace,placeins}
\usepackage{array,booktabs,tabularx}
\usepackage{xcolor}
\usepackage{microtype}
\usepackage{enumitem}
\usepackage{listings}
\usepackage[most]{tcolorbox}
\usepackage{titlesec}
\usepackage{fancyhdr}
\usepackage{hyperref}
\usepackage{bookmark}

\definecolor{ink}{HTML}{243041}
\definecolor{accent}{HTML}{38598C}
\definecolor{secondary}{HTML}{60768A}
\definecolor{codebg}{HTML}{F4F6F8}
\definecolor{codeborder}{HTML}{E3E8EE}
\definecolor{notebg}{HTML}{F3F7FC}

\hypersetup{
  colorlinks=true,
  linkcolor=accent,
  urlcolor=accent,
  citecolor=accent,
  pdfborder={0 0 0}
}
\urlstyle{same}

\setlength{\parindent}{0pt}
\setlength{\parskip}{0.38em plus 0.12em minus 0.05em}

\setlist[itemize]{
  topsep=0.35em,
  itemsep=0.13em,
  leftmargin=1.7em
}
\setlist[enumerate]{
  topsep=0.35em,
  itemsep=0.13em,
  leftmargin=1.8em
}

\setcounter{tocdepth}{2}
\setcounter{secnumdepth}{2}

\titleformat{\section}
  {\Large\bfseries\color{ink}}
  {\thesection}{0.7em}{}

\titleformat{\subsection}
  {\large\bfseries\color{accent}}
  {\thesubsection}{0.7em}{}

\titleformat{\subsubsection}
  {\normalsize\bfseries\color{ink}}
  {\thesubsubsection}{0.6em}{}

\titlespacing*{\section}
  {0pt}{2.2ex plus 0.7ex}{0.9ex}
\titlespacing*{\subsection}
  {0pt}{1.7ex plus 0.4ex}{0.6ex}
\titlespacing*{\subsubsection}
  {0pt}{1.3ex plus 0.3ex}{0.45ex}

\pagestyle{fancy}
\fancyhf{}

\fancyhead[L]{
  \footnotesize\sffamily\color{secondary}
  HF Deep RL Course \enspace/\enspace Unit 1
}
\fancyhead[R]{
  \footnotesize\sffamily\color{secondary}
  \nouppercase{\rightmark}
}
\fancyfoot[C]{
  \footnotesize\color{secondary}\thepage
}
\renewcommand{\headrulewidth}{0.25pt}

\captionsetup{
  font=small,
  labelfont={bf,color=secondary},
  textfont={color=ink},
  justification=centering,
  skip=4pt,
  hypcap=false
}

\lstdefinestyle{hfcode}{
  basicstyle=\ttfamily\footnotesize,
  backgroundcolor=\color{codebg},
  frame=single,
  rulecolor=\color{codeborder},
  numbers=none,
  showstringspaces=false,
  breaklines=true,
  breakatwhitespace=false,
  columns=fullflexible,
  keepspaces=true,
  xleftmargin=0pt,
  xrightmargin=0pt,
  aboveskip=0.6em,
  belowskip=0.6em,
  framexleftmargin=7pt,
  framexrightmargin=7pt,
  framextopmargin=5pt,
  framexbottommargin=5pt,
  upquote=true,
  tabsize=4,
}
\lstset{style=hfcode}

\newcommand{\CourseFigure}[2]{%
  \par\addvspace{0.55em}\noindent%
  \begin{minipage}{\linewidth}%
    \centering
    \includegraphics[
      width=0.86\linewidth,
      height=0.295\textheight,
      keepaspectratio
    ]{#1}\par
    \captionof{figure}{#2}%
  \end{minipage}%
  \par\addvspace{0.55em}%
}

\newtcolorbox{coursenote}{
  enhanced,
  breakable,
  colback=notebg,
  colframe=accent!30,
  boxrule=0.35pt,
  arc=1.2mm,
  left=9pt,
  right=9pt,
  top=6pt,
  bottom=6pt
}

\renewenvironment{quote}
  {\begin{coursenote}}
  {\end{coursenote}}

\emergencystretch=2.5em

\providecommand{\tightlist}{
  \setlength{\itemsep}{0pt}
  \setlength{\parskip}{0pt}
}

\providecommand{\passthrough}[1]{#1}

\Urlmuskip=0mu plus 6mu\relax

\graphicspath{{../assets/unit2/}}
% The shared preamble has Unit 1 in its header: override it here per-document.
\fancyhead[L]{\footnotesize\sffamily\color{secondary}HF Deep RL Course \enspace/\enspace Unit 2}
\title{\textbf{Deep Reinforcement Learning}\\[0.35em]\large Unit 2: Introduction to Q-Learning}
\author{Hugging Face Deep RL Course\\\small Faithful MDX-to-\LaTeX{} study edition}
\date{}
\begin{document}
\maketitle
\vspace{-0.8em}
\begin{center}\small\color{secondary}Original content: Hugging Face Deep RL Course.\end{center}
\tableofcontents
\vspace{1em}
\clearpage
\section{Introduction to Q-Learning}\label{introduction-to-q-learning}

\CourseFigure{thumbnail.jpg}{Unit 2 thumbnail \protect\href{https://huggingface.co/datasets/huggingface-deep-rl-course/course-images/resolve/main/en/unit3/thumbnail.jpg}{[original figure]}}

In the first unit of this class, we learned about Reinforcement Learning (RL), the RL process, and the different methods to solve an RL problem. We also \textbf{trained our first agents and uploaded them to the Hugging Face Hub.}

In this unit, we're going to \textbf{dive deeper into one of the Reinforcement Learning methods: value-based methods} and study our first RL algorithm: \textbf{Q-Learning.}

We'll also \textbf{implement our first RL agent from scratch}, a Q-Learning agent, and will train it in two environments:

\begin{enumerate}
\def\labelenumi{\arabic{enumi}.}
\tightlist
\item
  Frozen-Lake-v1 (non-slippery version): where our agent will need to \textbf{go from the starting state (S) to the goal state (G)} by walking only on frozen tiles (F) and avoiding holes (H).
\item
  An autonomous taxi: where our agent will need \textbf{to learn to navigate} a city to \textbf{transport its passengers from point A to point B.}
\end{enumerate}

\CourseFigure{envs.png}{envs \protect\href{https://huggingface.co/datasets/huggingface-deep-rl-course/course-images/resolve/main/en/unit3/envs.gif}{[original animation]}}

Concretely, we will:

\begin{itemize}
\tightlist
\item
  Learn about \textbf{value-based methods}.
\item
  Learn about the \textbf{differences between Monte Carlo and Temporal Difference Learning}.
\item
  Study and implement \textbf{our first RL algorithm}: Q-Learning.
\end{itemize}

This unit is \textbf{fundamental if you want to be able to work on Deep Q-Learning}: the first Deep RL algorithm that played Atari games and beat the human level on some of them (breakout, space invaders, etc).

So let's get started!

\section{What is RL? A short recap}\label{what-is-rl-a-short-recap}

In RL, we build an agent that can \textbf{make smart decisions}. For instance, an agent that \textbf{learns to play a video game.} Or a trading agent that \textbf{learns to maximize its benefits} by deciding on \textbf{what stocks to buy and when to sell.}

\CourseFigure{rl-process.jpg}{The agent-environment feedback loop \protect\href{https://huggingface.co/datasets/huggingface-deep-rl-course/course-images/resolve/main/en/unit3/rl-process.jpg}{[original figure]}}

To make intelligent decisions, our agent will learn from the environment by \textbf{interacting with it through trial and error} and receiving rewards (positive or negative) \textbf{as unique feedback.}

Its goal \textbf{is to maximize its expected cumulative reward} (because of the reward hypothesis).

\textbf{The agent's decision-making process is called the policy \(\pi\):} given a state, a policy will output an action or a probability distribution over actions. That is, given an observation of the environment, a policy will provide an action (or multiple probabilities for each action) that the agent should take.

\CourseFigure{policy.jpg}{A policy maps observations to actions \protect\href{https://huggingface.co/datasets/huggingface-deep-rl-course/course-images/resolve/main/en/unit3/policy.jpg}{[original figure]}}

\textbf{Our goal is to find an optimal policy \(\pi^*\) }, aka., a policy that leads to the best expected cumulative reward.

And to find this optimal policy (hence solving the RL problem), there \textbf{are two main types of RL methods}:

\begin{itemize}
\tightlist
\item
  \emph{Policy-based methods}: \textbf{Train the policy directly} to learn which action to take given a state.
\item
  \emph{Value-based methods}: \textbf{Train a value function} to learn \textbf{which state is more valuable} and use this value function \textbf{to take the action that leads to it.}
\end{itemize}

\CourseFigure{two-approaches.jpg}{The two fundamental approaches: direct policy learning and value-function learning \protect\href{https://huggingface.co/datasets/huggingface-deep-rl-course/course-images/resolve/main/en/unit3/two-approaches.jpg}{[original figure]}}

And in this unit, \textbf{we'll dive deeper into the value-based methods.}

\section{Two types of value-based methods}\label{two-types-of-value-based-methods}

In value-based methods, \textbf{we learn a value function} that \textbf{maps a state to the expected value of being at that state.}

\CourseFigure{vbm-1.jpg}{Overview of value-based reinforcement learning \protect\href{https://huggingface.co/datasets/huggingface-deep-rl-course/course-images/resolve/main/en/unit3/vbm-1.jpg}{[original figure]}}

The value of a state is the \textbf{expected discounted return} the agent can get if it \textbf{starts at that state and then acts according to our policy.}

\begin{quote}
\textbf{Course note.} But what does it mean to act according to our policy? After all, we don't have a policy in value-based methods since we train a value function and not a policy.
\end{quote}

Remember that the goal of an \textbf{RL agent is to have an optimal policy \(\pi^*\).}

To find the optimal policy, we learned about two different methods:

\begin{itemize}
\tightlist
\item
  \emph{Policy-based methods:} \textbf{Directly train the policy} to select what action to take given a state (or a probability distribution over actions at that state). In this case, we \textbf{don't have a value function.}
\end{itemize}

\CourseFigure{two-approaches-2.jpg}{Policy-based learning: actions predicted directly by a policy \protect\href{https://huggingface.co/datasets/huggingface-deep-rl-course/course-images/resolve/main/en/unit3/two-approaches-2.jpg}{[original figure]}}

The policy takes a state as input and outputs what action to take at that state (deterministic policy: a policy that output one action given a state, contrary to stochastic policy that output a probability distribution over actions).

And consequently, \textbf{we don't define by hand the behavior of our policy; it's the training that will define it.}

\begin{itemize}
\tightlist
\item
  \emph{Value-based methods:} \textbf{Indirectly, by training a value function} that outputs the value of a state or a state-action pair. Given this value function, our policy \textbf{will take an action.}
\end{itemize}

Since the policy is not trained/learned, \textbf{we need to specify its behavior.} For instance, if we want a policy that, given the value function, will take actions that always lead to the biggest reward, \textbf{we'll create a Greedy Policy.}

\CourseFigure{two-approaches-3.jpg}{Given a state, our action-value function (that we train) outputs the value of each action at that state. Then, our pre-defined Greedy Policy selects the action that will yield the highest value given a state or a state action pair. \protect\href{https://huggingface.co/datasets/huggingface-deep-rl-course/course-images/resolve/main/en/unit3/two-approaches-3.jpg}{[original figure]}}

Consequently, whatever method you use to solve your problem, \textbf{you will have a policy}. In the case of value-based methods, you don't train the policy: your policy \textbf{is just a simple pre-specified function} (for instance, the Greedy Policy) that uses the values given by the value-function to select its actions.

So the difference is:

\begin{itemize}
\tightlist
\item
  In policy-based training, \textbf{the optimal policy (denoted \(\pi^*\)) is found by training the policy directly.}
\item
  In value-based training, \textbf{finding an optimal value function (denoted Q* or V*, we'll study the difference below) leads to having an optimal policy.}
\end{itemize}

\CourseFigure{link-value-policy.jpg}{Recovering an optimal policy from the optimal action-value function \protect\href{https://huggingface.co/datasets/huggingface-deep-rl-course/course-images/resolve/main/en/unit3/link-value-policy.jpg}{[original figure]}}

In fact, most of the time, in value-based methods, you'll use \textbf{an Epsilon-Greedy Policy} that handles the exploration/exploitation trade-off; we'll talk about this when we talk about Q-Learning in the second part of this unit.

As we mentioned above, we have two types of value-based functions:

\subsection{The state-value function}\label{the-state-value-function}

We write the state value function under a policy \(\pi\) like this:

\CourseFigure{state-value-function-1.jpg}{Mathematical definition of the state-value function \protect\href{https://huggingface.co/datasets/huggingface-deep-rl-course/course-images/resolve/main/en/unit3/state-value-function-1.jpg}{[original figure]}}

For each state, the state-value function outputs the expected return if the agent \textbf{starts at that state} and then follows the policy forever afterward (for all future timesteps, if you prefer).

\CourseFigure{state-value-function-2.jpg}{If we take the state with value -7: it's the expected return starting at that state and taking actions according to our policy (greedy policy), so right, right, right, down, down, right, right. \protect\href{https://huggingface.co/datasets/huggingface-deep-rl-course/course-images/resolve/main/en/unit3/state-value-function-2.jpg}{[original figure]}}

\subsection{The action-value function}\label{the-action-value-function}

In the action-value function, for each state and action pair, the action-value function \textbf{outputs the expected return} if the agent starts in that state, takes that action, and then follows the policy forever after.

The value of taking action \(a\) in state \(s\) under a policy \(\pi\) is:

\CourseFigure{action-state-value-function-1.jpg}{Mathematical definition of the action-value function \protect\href{https://huggingface.co/datasets/huggingface-deep-rl-course/course-images/resolve/main/en/unit3/action-state-value-function-1.jpg}{[original figure]}}

\CourseFigure{action-state-value-function-2.jpg}{Action values for a chosen state and action \protect\href{https://huggingface.co/datasets/huggingface-deep-rl-course/course-images/resolve/main/en/unit3/action-state-value-function-2.jpg}{[original figure]}}

We see that the difference is:

\begin{itemize}
\tightlist
\item
  For the state-value function, we calculate \textbf{the value of a state \(S_t\)}
\item
  For the action-value function, we calculate \textbf{the value of the state-action pair ( \(S_t, A_t\) ) hence the value of taking that action at that state.}
\end{itemize}

\CourseFigure{two-types.jpg}{Note: We didn't fill all the state-action pairs for the example of Action-value function \protect\href{https://huggingface.co/datasets/huggingface-deep-rl-course/course-images/resolve/main/en/unit3/two-types.jpg}{[original figure]}}

In either case, whichever value function we choose (state-value or action-value function), \textbf{the returned value is the expected return.}

However, the problem is that \textbf{to calculate EACH value of a state or a state-action pair, we need to sum all the rewards an agent can get if it starts at that state.}

This can be a computationally expensive process, and that's \textbf{where the Bellman equation comes in to help us.}

\section{The Bellman Equation: simplify our value estimation}\label{the-bellman-equation-simplify-our-value-estimation}

The Bellman equation \textbf{simplifies our state value or state-action value calculation.}

\CourseFigure{bellman.jpg}{Bellman equation and recursive value estimation \protect\href{https://huggingface.co/datasets/huggingface-deep-rl-course/course-images/resolve/main/en/unit3/bellman.jpg}{[original figure]}}

With what we have learned so far, we know that if we calculate \(V(S_t)\) (the value of a state), we need to calculate the return starting at that state and then follow the policy forever after. \textbf{(The policy we defined in the following example is a Greedy Policy; for simplification, we don't discount the reward).}

So to calculate \(V(S_t)\), we need to calculate the sum of the expected rewards. Hence:

\CourseFigure{bellman2.jpg}{To calculate the value of State 1: the sum of rewards if the agent started in that state and then followed the greedy policy (taking actions that leads to the best states values) for all the time steps. \protect\href{https://huggingface.co/datasets/huggingface-deep-rl-course/course-images/resolve/main/en/unit3/bellman2.jpg}{[original figure]}}

Then, to calculate the \(V(S_{t+1})\), we need to calculate the return starting at that state \(S_{t+1}\).

\CourseFigure{bellman3.jpg}{To calculate the value of State 2: the sum of rewards if the agent started in that state, and then followed the policy for all the time steps. \protect\href{https://huggingface.co/datasets/huggingface-deep-rl-course/course-images/resolve/main/en/unit3/bellman3.jpg}{[original figure]}}

So you may have noticed, we're repeating the computation of the value of different states, which can be tedious if you need to do it for each state value or state-action value.

Instead of calculating the expected return for each state or each state-action pair, \textbf{we can use the Bellman equation.} (hint: if you know what Dynamic Programming is, this is very similar! if you don't know what it is, no worries!)

The Bellman equation is a recursive equation that works like this: instead of starting for each state from the beginning and calculating the return, we can consider the value of any state as:

\textbf{The immediate reward \(R_{t+1}\) + the discounted value of the state that follows ( \(\gamma * V(S_{t+1})\) ) .}

\CourseFigure{bellman4.jpg}{Bellman decomposition into immediate reward and discounted next-state value \protect\href{https://huggingface.co/datasets/huggingface-deep-rl-course/course-images/resolve/main/en/unit3/bellman4.jpg}{[original figure]}}

If we go back to our example, we can say that the value of State 1 is equal to the expected cumulative return if we start at that state.

\CourseFigure{bellman2.jpg}{Return accumulated from the first state in a trajectory \protect\href{https://huggingface.co/datasets/huggingface-deep-rl-course/course-images/resolve/main/en/unit3/bellman2.jpg}{[original figure]}}

To calculate the value of State 1: the sum of rewards \textbf{if the agent started in that state 1} and then followed the \textbf{policy for all the time steps.}

This is equivalent to \(V(S_{t})\) = Immediate reward \(R_{t+1}\) + Discounted value of the next state \(\gamma * V(S_{t+1})\)

\CourseFigure{bellman6.jpg}{For simplification, here we don’t discount so gamma = 1. \protect\href{https://huggingface.co/datasets/huggingface-deep-rl-course/course-images/resolve/main/en/unit3/bellman6.jpg}{[original figure]}}

In the interest of simplicity, here we don't discount, so gamma = 1. But you'll study an example with gamma = 0.99 in the Q-Learning section of this unit.

\begin{itemize}
\tightlist
\item
  The value of \(V(S_{t+1})\) = Immediate reward \(R_{t+2}\) + Discounted value of the next state ( \(gamma * V(S_{t+2})\) ).
\item
  And so on.
\end{itemize}

To recap, the idea of the Bellman equation is that instead of calculating each value as the sum of the expected return, \textbf{which is a long process}, we calculate the value as \textbf{the sum of immediate reward + the discounted value of the state that follows.}

Before going to the next section, think about the role of gamma in the Bellman equation. What happens if the value of gamma is very low (e.g.~0.1 or even 0)? What happens if the value is 1? What happens if the value is very high, such as a million?

\section{Monte Carlo vs Temporal Difference Learning}\label{monte-carlo-vs-temporal-difference-learning}

The last thing we need to discuss before diving into Q-Learning is the two learning strategies.

Remember that an RL agent \textbf{learns by interacting with its environment.} The idea is that \textbf{given the experience and the received reward, the agent will update its value function or policy.}

Monte Carlo and Temporal Difference Learning are two different \textbf{strategies on how to train our value function or our policy function.} Both of them \textbf{use experience to solve the RL problem.}

On one hand, Monte Carlo uses \textbf{an entire episode of experience before learning.} On the other hand, Temporal Difference uses \textbf{only a step ( \(S_t, A_t, R_{t+1}, S_{t+1}\) ) to learn.}

We'll explain both of them \textbf{using a value-based method example.}

\subsection{Monte Carlo: learning at the end of the episode}\label{monte-carlo-learning-at-the-end-of-the-episode}

Monte Carlo waits until the end of the episode, calculates \(G_t\) (return) and uses it as \textbf{a target for updating \(V(S_t)\).}

So it requires a \textbf{complete episode of interaction before updating our value function.}

\CourseFigure{monte-carlo-approach.jpg}{Monte Carlo value estimation from complete episode returns \protect\href{https://huggingface.co/datasets/huggingface-deep-rl-course/course-images/resolve/main/en/unit3/monte-carlo-approach.jpg}{[original figure]}}

If we take an example:

\CourseFigure{MC-2.jpg}{Monte Carlo learning: worked example (stage 2) \protect\href{https://huggingface.co/datasets/huggingface-deep-rl-course/course-images/resolve/main/en/unit3/MC-2.jpg}{[original figure]}}

\begin{itemize}
\item
  We always start the episode \textbf{at the same starting point.}
\item
  \textbf{The agent takes actions using the policy}. For instance, using an Epsilon Greedy Strategy, a policy that alternates between exploration (random actions) and exploitation.
\item
  We get \textbf{the reward and the next state.}
\item
  We terminate the episode if the cat eats the mouse or if the mouse moves \textgreater{} 10 steps.
\item
  At the end of the episode, \textbf{we have a list of State, Actions, Rewards, and Next States tuples}
\end{itemize}

For instance {[}{[}State tile 3 bottom, Go Left, +1, State tile 2 bottom{]}, {[}State tile 2 bottom, Go Left, +0, State tile 1 bottom{]}\ldots{]}

\begin{itemize}
\tightlist
\item
  \textbf{The agent will sum the total rewards \(G_t\)} (to see how well it did).
\item
  It will then \textbf{update \(V(s_t)\) based on the formula}
\end{itemize}

\CourseFigure{MC-3.jpg}{Monte Carlo learning: worked example (stage 3) \protect\href{https://huggingface.co/datasets/huggingface-deep-rl-course/course-images/resolve/main/en/unit3/MC-3.jpg}{[original figure]}}

\begin{itemize}
\tightlist
\item
  Then \textbf{start a new game with this new knowledge}
\end{itemize}

By running more and more episodes, \textbf{the agent will learn to play better and better.}

\CourseFigure{MC-3p.jpg}{Monte Carlo learning: worked example (stage 3) \protect\href{https://huggingface.co/datasets/huggingface-deep-rl-course/course-images/resolve/main/en/unit3/MC-3p.jpg}{[original figure]}}

For instance, if we train a state-value function using Monte Carlo:

\begin{itemize}
\tightlist
\item
  We initialize our value function \textbf{so that it returns 0 value for each state}
\item
  Our learning rate (lr) is 0.1 and our discount rate is 1 (= no discount)
\item
  Our mouse \textbf{explores the environment and takes random actions}
\end{itemize}

\CourseFigure{MC-4.jpg}{Monte Carlo learning: worked example (stage 4) \protect\href{https://huggingface.co/datasets/huggingface-deep-rl-course/course-images/resolve/main/en/unit3/MC-4.jpg}{[original figure]}}

\begin{itemize}
\tightlist
\item
  The mouse made more than 10 steps, so the episode ends .
\end{itemize}

\CourseFigure{MC-4p.jpg}{Monte Carlo learning: worked example (stage 4) \protect\href{https://huggingface.co/datasets/huggingface-deep-rl-course/course-images/resolve/main/en/unit3/MC-4p.jpg}{[original figure]}}

\begin{itemize}
\tightlist
\item
  We have a list of state, action, rewards, next\_state, \textbf{we need to calculate the return \(G{t=0}\)}
\end{itemize}

\(G_t = R_{t+1} + R_{t+2} + R_{t+3} ...\) (for simplicity, we don't discount the rewards)

\(G_0 = R_{1} + R_{2} + R_{3}\ldots\)

\(G_0 = 1 + 0 + 0 + 0 + 0 + 0 + 1 + 1 + 0 + 0\)

\(G_0 = 3\)

\begin{itemize}
\tightlist
\item
  We can now compute the \textbf{new} \(V(S_0)\):
\end{itemize}

\CourseFigure{MC-5.jpg}{Monte Carlo learning: worked example (stage 5) \protect\href{https://huggingface.co/datasets/huggingface-deep-rl-course/course-images/resolve/main/en/unit3/MC-5.jpg}{[original figure]}}

\(V(S_0) = V(S_0) + lr * [G_0 - V(S_0)]\)

\(V(S_0) = 0 + 0.1 * [3 - 0]\)

\(V(S_0) = 0.3\)

\CourseFigure{MC-5p.jpg}{Monte Carlo learning: worked example (stage 5) \protect\href{https://huggingface.co/datasets/huggingface-deep-rl-course/course-images/resolve/main/en/unit3/MC-5p.jpg}{[original figure]}}

\subsection{Temporal Difference Learning: learning at each step}\label{temporal-difference-learning-learning-at-each-step}

\textbf{Temporal Difference, on the other hand, waits for only one interaction (one step) \(S_{t+1}\)} to form a TD target and update \(V(S_t)\) using \(R_{t+1}\) and \(\gamma * V(S_{t+1})\).

The idea with \textbf{TD is to update the \(V(S_t)\) at each step.}

But because we didn't experience an entire episode, we don't have \(G_t\) (expected return). Instead, \textbf{we estimate \(G_t\) by adding \(R_{t+1}\) and the discounted value of the next state.}

This is called bootstrapping. It's called this \textbf{because TD bases its update in part on an existing estimate \(V(S_{t+1})\) and not a complete sample \(G_t\).}

\CourseFigure{TD-1.jpg}{Temporal Difference learning: worked example (stage 1) \protect\href{https://huggingface.co/datasets/huggingface-deep-rl-course/course-images/resolve/main/en/unit3/TD-1.jpg}{[original figure]}}

This method is called TD(0) or \textbf{one-step TD (update the value function after any individual step).}

\CourseFigure{TD-1p.jpg}{Temporal Difference learning: worked example (stage 1) \protect\href{https://huggingface.co/datasets/huggingface-deep-rl-course/course-images/resolve/main/en/unit3/TD-1p.jpg}{[original figure]}}

If we take the same example,

\CourseFigure{TD-2.jpg}{Temporal Difference learning: worked example (stage 2) \protect\href{https://huggingface.co/datasets/huggingface-deep-rl-course/course-images/resolve/main/en/unit3/TD-2.jpg}{[original figure]}}

\begin{itemize}
\tightlist
\item
  We initialize our value function so that it returns 0 value for each state.
\item
  Our learning rate (lr) is 0.1, and our discount rate is 1 (no discount).
\item
  Our mouse begins to explore the environment and takes a random action: \textbf{going to the left}
\item
  It gets a reward \(R_{t+1} = 1\) since \textbf{it eats a piece of cheese}
\end{itemize}

\CourseFigure{TD-2p.jpg}{Temporal Difference learning: worked example (stage 2) \protect\href{https://huggingface.co/datasets/huggingface-deep-rl-course/course-images/resolve/main/en/unit3/TD-2p.jpg}{[original figure]}}

\CourseFigure{TD-3.jpg}{Temporal Difference learning: worked example (stage 3) \protect\href{https://huggingface.co/datasets/huggingface-deep-rl-course/course-images/resolve/main/en/unit3/TD-3.jpg}{[original figure]}}

We can now update \(V(S_0)\):

New \(V(S_0) = V(S_0) + lr * [R_1 + \gamma * V(S_1) - V(S_0)]\)

New \(V(S_0) = 0 + 0.1 * [1 + 1 * 0-0]\)

New \(V(S_0) = 0.1\)

So we just updated our value function for State 0.

Now we \textbf{continue to interact with this environment with our updated value function.}

\CourseFigure{TD-3p.jpg}{Temporal Difference learning: worked example (stage 3) \protect\href{https://huggingface.co/datasets/huggingface-deep-rl-course/course-images/resolve/main/en/unit3/TD-3p.jpg}{[original figure]}}

To summarize:

\begin{itemize}
\tightlist
\item
  With \emph{Monte Carlo}, we update the value function from a complete episode, and so we \textbf{use the actual accurate discounted return of this episode.}
\item
  With \emph{TD Learning}, we update the value function from a step, and we replace \(G_t\), which we don't know, with \textbf{an estimated return called the TD target.}
\end{itemize}

\CourseFigure{Summary.jpg}{Monte Carlo versus temporal-difference learning: comparison \protect\href{https://huggingface.co/datasets/huggingface-deep-rl-course/course-images/resolve/main/en/unit3/Summary.jpg}{[original figure]}}

\section{Mid-way Recap}\label{mid-way-recap}

Before diving into Q-Learning, let's summarize what we've just learned.

We have two types of value-based functions:

\begin{itemize}
\tightlist
\item
  State-value function: outputs the expected return if \textbf{the agent starts at a given state and acts according to the policy forever after.}
\item
  Action-value function: outputs the expected return if \textbf{the agent starts in a given state, takes a given action at that state} and then acts accordingly to the policy forever after.
\item
  In value-based methods, rather than learning the policy, \textbf{we define the policy by hand} and we learn a value function. If we have an optimal value function, we \textbf{will have an optimal policy.}
\end{itemize}

There are two types of methods to update the value function:

\begin{itemize}
\tightlist
\item
  With \emph{the Monte Carlo method}, we update the value function from a complete episode, and so we \textbf{use the actual discounted return of this episode.}
\item
  With \emph{the TD Learning method,} we update the value function from a step, replacing the unknown \(G_t\) with \textbf{an estimated return called the TD target.}
\end{itemize}

\CourseFigure{summary-learning-mtds.jpg}{Summary \protect\href{https://huggingface.co/datasets/huggingface-deep-rl-course/course-images/resolve/main/en/unit3/summary-learning-mtds.jpg}{[original figure]}}

\section{Mid-way Quiz}\label{mid-way-quiz}

The best way to learn and \href{https://www.coursera.org/lecture/learning-how-to-learn/illusions-of-competence-BuFzf}{to avoid the illusion of competence} \textbf{is to test yourself.} This will help you to find \textbf{where you need to reinforce your knowledge}.

\subsubsection{Q1: What are the two main approaches to find optimal policy?}\label{q1-what-are-the-two-main-approaches-to-find-optimal-policy}

\textbf{Select all applicable answers.}

\begin{itemize}
\tightlist
\item
  \textbf{A.} Policy-based methods
\item
  \textbf{B.} Random-based methods
\item
  \textbf{C.} Value-based methods
\item
  \textbf{D.} Evolution-strategies methods
\end{itemize}

\subsubsection{Q2: What is the Bellman Equation?}\label{q2-what-is-the-bellman-equation}

\subsubsection{Q3: Define each part of the Bellman Equation}\label{q3-define-each-part-of-the-bellman-equation}

\CourseFigure{bellman4-quiz.jpg}{Bellman equation quiz \protect\href{https://huggingface.co/datasets/huggingface-deep-rl-course/course-images/resolve/main/en/unit3/bellman4-quiz.jpg}{[original figure]}}

\subsubsection{Q4: What is the difference between Monte Carlo and Temporal Difference learning methods?}\label{q4-what-is-the-difference-between-monte-carlo-and-temporal-difference-learning-methods}

\textbf{Select all applicable answers.}

\begin{itemize}
\tightlist
\item
  \textbf{A.} With Monte Carlo methods, we update the value function from a complete episode
\item
  \textbf{B.} With Monte Carlo methods, we update the value function from a step
\item
  \textbf{C.} With TD learning methods, we update the value function from a complete episode
\item
  \textbf{D.} With TD learning methods, we update the value function from a step
\end{itemize}

\subsubsection{Q5: Define each part of Temporal Difference learning formula}\label{q5-define-each-part-of-temporal-difference-learning-formula}

\CourseFigure{td-ex.jpg}{TD Learning exercise \protect\href{https://huggingface.co/datasets/huggingface-deep-rl-course/course-images/resolve/main/en/unit3/td-ex.jpg}{[original figure]}}

\subsubsection{Q6: Define each part of Monte Carlo learning formula}\label{q6-define-each-part-of-monte-carlo-learning-formula}

\CourseFigure{mc-ex.jpg}{MC Learning exercise \protect\href{https://huggingface.co/datasets/huggingface-deep-rl-course/course-images/resolve/main/en/unit3/mc-ex.jpg}{[original figure]}}

\subsection{Answer Key and Explanations}\label{answer-key-and-explanations}

\subsubsection{Q1: What are the two main approaches to find optimal policy?}\label{q1-what-are-the-two-main-approaches-to-find-optimal-policy-1}

\textbf{Correct choice(s): A, C}

\begin{itemize}
\item
  \textbf{Option A (correct):} With Policy-Based methods, we train the policy directly to learn which action to take given a state.
\item
  \textbf{Option C (correct):} With value-based methods, we train a value function to learn which state is more valuable and use this value function to take the action that leads to it.
\end{itemize}

\subsubsection{Q2: What is the Bellman Equation?}\label{q2-what-is-the-bellman-equation-1}

\textbf{The Bellman equation is a recursive equation} that works like this: instead of starting for each state from the beginning and calculating the return, we can consider the value of any state as:

Rt+1 + gamma * V(St+1)

The immediate reward + the discounted value of the state that follows

\subsubsection{Q3: Define each part of the Bellman Equation}\label{q3-define-each-part-of-the-bellman-equation-1}

\CourseFigure{bellman4.jpg}{Bellman decomposition into immediate reward and discounted next-state value \protect\href{https://huggingface.co/datasets/huggingface-deep-rl-course/course-images/resolve/main/en/unit3/bellman4.jpg}{[original figure]}}

\subsubsection{Q4: What is the difference between Monte Carlo and Temporal Difference learning methods?}\label{q4-what-is-the-difference-between-monte-carlo-and-temporal-difference-learning-methods-1}

\textbf{Correct choice(s): A, D}

\subsubsection{Q5: Define each part of Temporal Difference learning formula}\label{q5-define-each-part-of-temporal-difference-learning-formula-1}

\CourseFigure{TD-1.jpg}{Temporal Difference learning: worked example (stage 1) \protect\href{https://huggingface.co/datasets/huggingface-deep-rl-course/course-images/resolve/main/en/unit3/TD-1.jpg}{[original figure]}}

\subsubsection{Q6: Define each part of Monte Carlo learning formula}\label{q6-define-each-part-of-monte-carlo-learning-formula-1}

\CourseFigure{monte-carlo-approach.jpg}{Monte Carlo value estimation from complete episode returns \protect\href{https://huggingface.co/datasets/huggingface-deep-rl-course/course-images/resolve/main/en/unit3/monte-carlo-approach.jpg}{[original figure]}}

Congrats on finishing this Quiz , if you missed some elements, take time to read again the previous sections to reinforce () your knowledge.

\section{Introducing Q-Learning}\label{introducing-q-learning}

\subsection{What is Q-Learning?}\label{what-is-q-learning}

Q-Learning is an \textbf{off-policy value-based method that uses a TD approach to train its action-value function:}

\begin{itemize}
\tightlist
\item
  \emph{Off-policy}: we'll talk about that at the end of this unit.
\item
  \emph{Value-based method}: finds the optimal policy indirectly by training a value or action-value function that will tell us \textbf{the value of each state or each state-action pair.}
\item
  \emph{TD approach:} \textbf{updates its action-value function at each step instead of at the end of the episode.}
\end{itemize}

\textbf{Q-Learning is the algorithm we use to train our Q-function}, an \textbf{action-value function} that determines the value of being at a particular state and taking a specific action at that state.

\CourseFigure{Q-function.jpg}{Given a state and action, our Q Function outputs a state-action value (also called Q-value) \protect\href{https://huggingface.co/datasets/huggingface-deep-rl-course/course-images/resolve/main/en/unit3/Q-function.jpg}{[original figure]}}

The \textbf{Q comes from ``the Quality'' (the value) of that action at that state.}

Let's recap the difference between value and reward:

\begin{itemize}
\tightlist
\item
  The \emph{value of a state}, or a \emph{state-action pair} is the expected cumulative reward our agent gets if it starts at this state (or state-action pair) and then acts accordingly to its policy.
\item
  The \emph{reward} is the \textbf{feedback it gets from the environment} after performing an action at a state.
\end{itemize}

Internally, our Q-function is encoded by \textbf{a Q-table, a table where each cell corresponds to a state-action pair value.} Think of this Q-table as \textbf{the memory or cheat sheet of our Q-function.}

Let's go through an example of a maze.

\CourseFigure{Maze-1.jpg}{Q-table illustrated on a mouse-and-cheese maze (Maze-1) \protect\href{https://huggingface.co/datasets/huggingface-deep-rl-course/course-images/resolve/main/en/unit3/Maze-1.jpg}{[original figure]}}

The Q-table is initialized. That's why all values are = 0. This table \textbf{contains, for each state and action, the corresponding state-action values.} For this simple example, the state is only defined by the position of the mouse. Therefore, we have 2*3 rows in our Q-table, one row for each possible position of the mouse. In more complex scenarios, the state could contain more information than the position of the actor.

\CourseFigure{Maze-2.jpg}{Q-table illustrated on a mouse-and-cheese maze (Maze-2) \protect\href{https://huggingface.co/datasets/huggingface-deep-rl-course/course-images/resolve/main/en/unit3/Maze-2.jpg}{[original figure]}}

Here we see that the \textbf{state-action value of the initial state and going up is 0:}

\CourseFigure{Maze-3.jpg}{Q-table illustrated on a mouse-and-cheese maze (Maze-3) \protect\href{https://huggingface.co/datasets/huggingface-deep-rl-course/course-images/resolve/main/en/unit3/Maze-3.jpg}{[original figure]}}

So: the Q-function uses a Q-table \textbf{that has the value of each state-action pair.} Given a state and action, \textbf{our Q-function will search inside its Q-table to output the value.}

\CourseFigure{Q-function-2.jpg}{Q-table lookup: obtaining the value of a state-action pair \protect\href{https://huggingface.co/datasets/huggingface-deep-rl-course/course-images/resolve/main/en/unit3/Q-function-2.jpg}{[original figure]}}

If we recap, \emph{Q-Learning} \textbf{is the RL algorithm that:}

\begin{itemize}
\tightlist
\item
  Trains a \emph{Q-function} (an \textbf{action-value function}), which internally is a \textbf{Q-table that contains all the state-action pair values.}
\item
  Given a state and action, our Q-function \textbf{will search its Q-table for the corresponding value.}
\item
  When the training is done, \textbf{we have an optimal Q-function, which means we have optimal Q-table.}
\item
  And if we \textbf{have an optimal Q-function}, we \textbf{have an optimal policy} since we \textbf{know the best action to take at each state.}
\end{itemize}

\CourseFigure{link-value-policy.jpg}{Recovering an optimal policy from the optimal action-value function \protect\href{https://huggingface.co/datasets/huggingface-deep-rl-course/course-images/resolve/main/en/unit3/link-value-policy.jpg}{[original figure]}}

In the beginning, \textbf{our Q-table is useless since it gives arbitrary values for each state-action pair} (most of the time, we initialize the Q-table to 0). As the agent \textbf{explores the environment and we update the Q-table, it will give us a better and better approximation} to the optimal policy.

\CourseFigure{Q-learning-1.jpg}{We see here that with the training, our Q-table is better since, thanks to it, we can know the value of each state-action pair. \protect\href{https://huggingface.co/datasets/huggingface-deep-rl-course/course-images/resolve/main/en/unit3/Q-learning-1.jpg}{[original figure]}}

Now that we understand what Q-Learning, Q-functions, and Q-tables are, \textbf{let's dive deeper into the Q-Learning algorithm}.

\subsection{The Q-Learning algorithm}\label{the-q-learning-algorithm}

This is the Q-Learning pseudocode; let's study each part and \textbf{see how it works with a simple example before implementing it.} Don't be intimidated by it, it's simpler than it looks! We'll go over each step.

\CourseFigure{Q-learning-2.jpg}{Q-learning: reference algorithm and update pseudocode \protect\href{https://huggingface.co/datasets/huggingface-deep-rl-course/course-images/resolve/main/en/unit3/Q-learning-2.jpg}{[original figure]}}

\subsubsection{Step 1: We initialize the Q-table}\label{step-1-we-initialize-the-q-table}

\CourseFigure{Q-learning-3.jpg}{Q-learning: choose action using epsilon-greedy policy \protect\href{https://huggingface.co/datasets/huggingface-deep-rl-course/course-images/resolve/main/en/unit3/Q-learning-3.jpg}{[original figure]}}

We need to initialize the Q-table for each state-action pair. \textbf{Most of the time, we initialize with values of 0.}

\subsubsection{Step 2: Choose an action using the epsilon-greedy strategy}\label{step-2-choose-an-action-using-the-epsilon-greedy-strategy}

\CourseFigure{Q-learning-4.jpg}{Epsilon-greedy strategy: exploration and exploitation \protect\href{https://huggingface.co/datasets/huggingface-deep-rl-course/course-images/resolve/main/en/unit3/Q-learning-4.jpg}{[original figure]}}

The epsilon-greedy strategy is a policy that handles the exploration/exploitation trade-off.

The idea is that, with an initial value of \(\varepsilon\) = 1.0:

\begin{itemize}
\tightlist
\item
  \emph{With probability 1 --- \(\varepsilon\)} : we do \textbf{exploitation} (aka our agent selects the action with the highest state-action pair value).
\item
  With probability \(\varepsilon\): \textbf{we do exploration} (trying random action).
\end{itemize}

At the beginning of the training, \textbf{the probability of doing exploration will be huge since \(\varepsilon\) is very high, so most of the time, we'll explore.} But as the training goes on, and consequently our \textbf{Q-table gets better and better in its estimations, we progressively reduce the epsilon value} since we will need less and less exploration and more exploitation.

\CourseFigure{Q-learning-5.jpg}{Q-learning: perform action and observe transition \protect\href{https://huggingface.co/datasets/huggingface-deep-rl-course/course-images/resolve/main/en/unit3/Q-learning-5.jpg}{[original figure]}}

\subsubsection{Step 3: Perform action At, get reward Rt+1 and next state St+1}\label{step-3-perform-action-at-get-reward-rt1-and-next-state-st1}

\CourseFigure{Q-learning-6.jpg}{Q-learning: update the action value using a TD target \protect\href{https://huggingface.co/datasets/huggingface-deep-rl-course/course-images/resolve/main/en/unit3/Q-learning-6.jpg}{[original figure]}}

\subsubsection{Step 4: Update Q(St, At)}\label{step-4-update-qst-at}

Remember that in TD Learning, we update our policy or value function (depending on the RL method we choose) \textbf{after one step of the interaction.}

To produce our TD target, \textbf{we used the immediate reward \(R_{t+1}\) plus the discounted value of the next state}, computed by finding the action that maximizes the current Q-function at the next state. (We call that bootstrap).

\CourseFigure{Q-learning-7.jpg}{Q-learning: repeat updates over multiple episodes \protect\href{https://huggingface.co/datasets/huggingface-deep-rl-course/course-images/resolve/main/en/unit3/Q-learning-7.jpg}{[original figure]}}

Therefore, our \(Q(S_t, A_t)\) \textbf{update formula goes like this:}

\CourseFigure{Q-learning-8.jpg}{Q-learning: complete algorithm overview \protect\href{https://huggingface.co/datasets/huggingface-deep-rl-course/course-images/resolve/main/en/unit3/Q-learning-8.jpg}{[original figure]}}

This means that to update our \(Q(S_t, A_t)\):

\begin{itemize}
\tightlist
\item
  We need \(S_t, A_t, R_{t+1}, S_{t+1}\).
\item
  To update our Q-value at a given state-action pair, we use the TD target.
\end{itemize}

How do we form the TD target?

\begin{enumerate}
\def\labelenumi{\arabic{enumi}.}
\tightlist
\item
  We obtain the reward \(R_{t+1}\) after taking the action \(A_t\).
\item
  To get the \textbf{best state-action pair value} for the next state, we use a greedy policy to select the next best action. Note that this is not an epsilon-greedy policy, this will always take the action with the highest state-action value.
\end{enumerate}

Then when the update of this Q-value is done, we start in a new state and select our action \textbf{using a epsilon-greedy policy again.}

\textbf{This is why we say that Q Learning is an off-policy algorithm.}

\subsection{Off-policy vs On-policy}\label{off-policy-vs-on-policy}

The difference is subtle:

\begin{itemize}
\tightlist
\item
  \emph{Off-policy}: using \textbf{a different policy for acting (inference) and updating (training).}
\end{itemize}

For instance, with Q-Learning, the epsilon-greedy policy (acting policy), is different from the greedy policy that is \textbf{used to select the best next-state action value to update our Q-value (updating policy).}

\CourseFigure{off-on-1.jpg}{Off-on policy \protect\href{https://huggingface.co/datasets/huggingface-deep-rl-course/course-images/resolve/main/en/unit3/off-on-1.jpg}{[original figure]}}

Is different from the policy we use during the training part:

\CourseFigure{off-on-2.jpg}{Off-on policy \protect\href{https://huggingface.co/datasets/huggingface-deep-rl-course/course-images/resolve/main/en/unit3/off-on-2.jpg}{[original figure]}}

\begin{itemize}
\tightlist
\item
  \emph{On-policy:} using the \textbf{same policy for acting and updating.}
\end{itemize}

For instance, with Sarsa, another value-based algorithm, \textbf{the epsilon-greedy policy selects the next state-action pair, not a greedy policy.}

\CourseFigure{off-on-3.jpg}{Off-on policy \protect\href{https://huggingface.co/datasets/huggingface-deep-rl-course/course-images/resolve/main/en/unit3/off-on-3.jpg}{[original figure]}}

\CourseFigure{off-on-4.jpg}{Comparison of off-policy and on-policy methods \protect\href{https://huggingface.co/datasets/huggingface-deep-rl-course/course-images/resolve/main/en/unit3/off-on-4.jpg}{[original figure]}}

\section{A Q-Learning example}\label{a-q-learning-example}

To better understand Q-Learning, let's take a simple example:

\CourseFigure{Maze-Example-2.jpg}{The mouse-and-cheese maze with reward and termination rules \protect\href{https://huggingface.co/datasets/huggingface-deep-rl-course/course-images/resolve/main/en/unit3/Maze-Example-2.jpg}{[original figure]}}

\begin{itemize}
\tightlist
\item
  You're a mouse in this tiny maze. You always \textbf{start at the same starting point.}
\item
  The goal is \textbf{to eat the big pile of cheese at the bottom right-hand corner} and avoid the poison. After all, who doesn't like cheese?
\item
  The episode ends if we eat the poison, \textbf{eat the big pile of cheese}, or if we take more than five steps.
\item
  The learning rate is 0.1
\item
  The discount rate (gamma) is 0.99
\end{itemize}

\CourseFigure{q-ex-1.jpg}{Q-learning example: transition and Q-table update (step 1) \protect\href{https://huggingface.co/datasets/huggingface-deep-rl-course/course-images/resolve/main/en/unit3/q-ex-1.jpg}{[original figure]}}

The reward function goes like this:

\begin{itemize}
\tightlist
\item
  \textbf{+0:} Going to a state with no cheese in it.
\item
  \textbf{+1:} Going to a state with a small cheese in it.
\item
  \textbf{+10:} Going to the state with the big pile of cheese.
\item
  \textbf{-10:} Going to the state with the poison and thus dying.
\item
  \textbf{+0} If we take more than five steps.
\end{itemize}

\CourseFigure{q-ex-2.jpg}{Q-learning example: transition and Q-table update (step 2) \protect\href{https://huggingface.co/datasets/huggingface-deep-rl-course/course-images/resolve/main/en/unit3/q-ex-2.jpg}{[original figure]}}

To train our agent to have an optimal policy (so a policy that goes right, right, down), \textbf{we will use the Q-Learning algorithm}.

\subsection{Step 1: Initialize the Q-table}\label{step-1-initialize-the-q-table}

\CourseFigure{Example-1.jpg}{Maze-Example \protect\href{https://huggingface.co/datasets/huggingface-deep-rl-course/course-images/resolve/main/en/unit3/Example-1.jpg}{[original figure]}}

So, for now, \textbf{our Q-table is useless}; we need \textbf{to train our Q-function using the Q-Learning algorithm.}

Let's do it for 2 training timesteps:

Training timestep 1:

\subsection{Step 2: Choose an action using the Epsilon Greedy Strategy}\label{step-2-choose-an-action-using-the-epsilon-greedy-strategy-1}

Because epsilon is big (= 1.0), I take a random action. In this case, I go right.

\CourseFigure{q-ex-3.jpg}{Q-learning example: transition and Q-table update (step 3) \protect\href{https://huggingface.co/datasets/huggingface-deep-rl-course/course-images/resolve/main/en/unit3/q-ex-3.jpg}{[original figure]}}

\subsection{Step 3: Perform action At, get Rt+1 and St+1}\label{step-3-perform-action-at-get-rt1-and-st1}

By going right, I get a small cheese, so \(R_{t+1} = 1\) and I'm in a new state.

\CourseFigure{q-ex-4.jpg}{Q-learning example: transition and Q-table update (step 4) \protect\href{https://huggingface.co/datasets/huggingface-deep-rl-course/course-images/resolve/main/en/unit3/q-ex-4.jpg}{[original figure]}}

\subsection{Step 4: Update Q(St, At)}\label{step-4-update-qst-at-1}

We can now update \(Q(S_t, A_t)\) using our formula.

\CourseFigure{q-ex-5.jpg}{Q-learning example: transition and Q-table update (step 5) \protect\href{https://huggingface.co/datasets/huggingface-deep-rl-course/course-images/resolve/main/en/unit3/q-ex-5.jpg}{[original figure]}}

\CourseFigure{Example-4.jpg}{Maze-Example \protect\href{https://huggingface.co/datasets/huggingface-deep-rl-course/course-images/resolve/main/en/unit3/Example-4.jpg}{[original figure]}}

Training timestep 2:

\subsection{Step 2: Choose an action using the Epsilon Greedy Strategy}\label{step-2-choose-an-action-using-the-epsilon-greedy-strategy-2}

\textbf{I take a random action again, since epsilon=0.99 is big}. (Notice we decay epsilon a little bit because, as the training progress, we want less and less exploration).

I took the action `down'. \textbf{This is not a good action since it leads me to the poison.}

\CourseFigure{q-ex-6.jpg}{Q-learning example: transition and Q-table update (step 6) \protect\href{https://huggingface.co/datasets/huggingface-deep-rl-course/course-images/resolve/main/en/unit3/q-ex-6.jpg}{[original figure]}}

\subsection{Step 3: Perform action At, get Rt+1 and St+1}\label{step-3-perform-action-at-get-rt1-and-st1-1}

Because I ate poison, \textbf{I get \(R_{t+1} = -10\), and I die.}

\CourseFigure{q-ex-7.jpg}{Q-learning example: transition and Q-table update (step 7) \protect\href{https://huggingface.co/datasets/huggingface-deep-rl-course/course-images/resolve/main/en/unit3/q-ex-7.jpg}{[original figure]}}

\subsection{Step 4: Update Q(St, At)}\label{step-4-update-qst-at-2}

\CourseFigure{q-ex-8.jpg}{Q-learning example: transition and Q-table update (step 8) \protect\href{https://huggingface.co/datasets/huggingface-deep-rl-course/course-images/resolve/main/en/unit3/q-ex-8.jpg}{[original figure]}}

Because we're dead, we start a new episode. But what we see here is that, \textbf{with two explorations steps, my agent became smarter.}

As we continue exploring and exploiting the environment and updating Q-values using the TD target, the \textbf{Q-table will give us a better and better approximation. At the end of the training, we'll get an estimate of the optimal Q-function.}

\section{Q-Learning Recap}\label{q-learning-recap}

\emph{Q-Learning} \textbf{is the RL algorithm that} :

\begin{itemize}
\item
  Trains a \emph{Q-function}, an \textbf{action-value function} encoded, in internal memory, by a \emph{Q-table} \textbf{containing all the state-action pair values.}
\item
  Given a state and action, our Q-function \textbf{will search its Q-table for the corresponding value.}
\end{itemize}

\CourseFigure{Q-function-2.jpg}{Q-table lookup: obtaining the value of a state-action pair \protect\href{https://huggingface.co/datasets/huggingface-deep-rl-course/course-images/resolve/main/en/unit3/Q-function-2.jpg}{[original figure]}}

\begin{itemize}
\item
  When the training is done, \textbf{we have an optimal Q-function, or, equivalently, an optimal Q-table.}
\item
  And if we \textbf{have an optimal Q-function}, we
\end{itemize}

have an optimal policy, since we \textbf{know, for each state, the best action to take.}

\CourseFigure{link-value-policy.jpg}{Recovering an optimal policy from the optimal action-value function \protect\href{https://huggingface.co/datasets/huggingface-deep-rl-course/course-images/resolve/main/en/unit3/link-value-policy.jpg}{[original figure]}}

But, in the beginning, our \textbf{Q-table is useless since it gives arbitrary values for each state-action pair (most of the time we initialize the Q-table to 0 values)}. But, as we explore the environment and update our Q-table it will give us a better and better approximation.

\CourseFigure{q-learning.jpeg}{Learning a Q-table from zero initialization toward a useful policy \protect\href{https://huggingface.co/datasets/huggingface-deep-rl-course/course-images/resolve/main/en/notebooks/unit2/q-learning.jpeg}{[original figure]}}

This is the Q-Learning pseudocode:

\CourseFigure{Q-learning-2.jpg}{Q-learning: reference algorithm and update pseudocode \protect\href{https://huggingface.co/datasets/huggingface-deep-rl-course/course-images/resolve/main/en/unit3/Q-learning-2.jpg}{[original figure]}}

\section{Hands-on}\label{hands-on}

\textbf{Official practical notebook:} \url{https://colab.research.google.com/github/huggingface/deep-rl-class/blob/main/notebooks/unit2/unit2.ipynb}

\textbf{Course help:} \url{http://hf.co/join/discord}

Now that we studied the Q-Learning algorithm, let's implement it from scratch and train our Q-Learning agent in two environments:

\begin{enumerate}
\def\labelenumi{\arabic{enumi}.}
\tightlist
\item
  \href{https://gymnasium.farama.org/environments/toy_text/frozen_lake/}{Frozen-Lake-v1 (non-slippery and slippery version)} : where our agent will need to \textbf{go from the starting state (S) to the goal state (G)} by walking only on frozen tiles (F) and avoiding holes (H).
\item
  \href{https://gymnasium.farama.org/environments/toy_text/taxi/}{An autonomous taxi} will need \textbf{to learn to navigate} a city to \textbf{transport its passengers from point A to point B.}
\end{enumerate}

\CourseFigure{envs.png}{envs \protect\href{https://huggingface.co/datasets/huggingface-deep-rl-course/course-images/resolve/main/en/unit3/envs.gif}{[original animation]}}

Thanks to a \href{https://huggingface.co/spaces/huggingface-projects/Deep-Reinforcement-Learning-Leaderboard}{leaderboard}, you'll be able to compare your results with other classmates and exchange the best practices to improve your agent's scores. Who will win the challenge for Unit 2?

To validate this hands-on for the \href{https://huggingface.co/deep-rl-course/en/unit0/introduction\#certification-process}{certification process}, you need to push your trained Taxi model to the Hub and \textbf{get a result of \textgreater= 4.5}.

To find your result, go to the \href{https://huggingface.co/spaces/huggingface-projects/Deep-Reinforcement-Learning-Leaderboard}{leaderboard} and find your model, \textbf{the result = mean\_reward - std of reward}

For more information about the certification process, check this section \url{https://huggingface.co/deep-rl-course/en/unit0/introduction\#certification-process}

And you can check your progress here \url{https://huggingface.co/spaces/ThomasSimonini/Check-my-progress-Deep-RL-Course}

\textbf{To start the hands-on click on the Open In Colab button} :

\href{https://colab.research.google.com/github/huggingface/deep-rl-class/blob/master/notebooks/unit2/unit2.ipynb}{Open original notebook in Google Colab}

We strongly \textbf{recommend students use Google Colab for the hands-on exercises} instead of running them on their personal computers.

By using Google Colab, \textbf{you can focus on learning and experimenting without worrying about the technical aspects} of setting up your environments.

\section{Unit 2: Q-Learning with FrozenLake-v1 and Taxi-v3}\label{unit-2-q-learning-with-frozenlake-v1-and-taxi-v3}

\CourseFigure{thumbnail.jpg}{Unit 2 Thumbnail \protect\href{https://huggingface.co/datasets/huggingface-deep-rl-course/course-images/resolve/main/en/unit3/thumbnail.jpg}{[original figure]}}

In this notebook, \textbf{you'll code your first Reinforcement Learning agent from scratch} to play FrozenLake using Q-Learning, share it with the community, and experiment with different configurations.

Here is an example of what \textbf{you will achieve in just a couple of minutes.}

\CourseFigure{envs.png}{envs \protect\href{https://huggingface.co/datasets/huggingface-deep-rl-course/course-images/resolve/main/en/unit3/envs.gif}{[original animation]}}

\subsubsection{Environments:}\label{environments}

\begin{itemize}
\tightlist
\item
  \href{https://gymnasium.farama.org/environments/toy_text/frozen_lake/}{FrozenLake-v1}
\item
  \href{https://gymnasium.farama.org/environments/toy_text/taxi/}{Taxi-v3}
\end{itemize}

\subsubsection{RL-Library:}\label{rl-library}

\begin{itemize}
\tightlist
\item
  Python and NumPy
\item
  \href{https://gymnasium.farama.org/}{Gymnasium}
\end{itemize}

We're constantly trying to improve our tutorials, so \textbf{if you find some issues in this notebook}, please \href{https://github.com/huggingface/deep-rl-class/issues}{open an issue on the GitHub Repo}.

\subsection{Objectives of this notebook}\label{objectives-of-this-notebook}

At the end of the notebook, you will:

\begin{itemize}
\tightlist
\item
  Be able to use \textbf{Gymnasium}, the environment library.
\item
  Be able to code a Q-Learning agent from scratch.
\item
  Be able to \textbf{push your trained agent and the code to the Hub} with a nice video replay and an evaluation score .
\end{itemize}

\subsection{This notebook is from the Deep Reinforcement Learning Course}\label{this-notebook-is-from-the-deep-reinforcement-learning-course}

\CourseFigure{deep-rl-course-illustration.jpg}{Deep RL Course illustration \protect\href{https://huggingface.co/datasets/huggingface-deep-rl-course/course-images/resolve/main/en/notebooks/deep-rl-course-illustration.jpg}{[original figure]}}

In this free course, you will:

\begin{itemize}
\tightlist
\item
  Study Deep Reinforcement Learning in \textbf{theory and practice}.
\item
  Learn to \textbf{use famous Deep RL libraries} such as Stable Baselines3, RL Baselines3 Zoo, CleanRL and Sample Factory 2.0.
\item
  Train \textbf{agents in unique environments}
\end{itemize}

And more check the syllabus \url{https://simoninithomas.github.io/deep-rl-course}

Don't forget to \textbf{\href{http://eepurl.com/ic5ZUD}{sign up to the course}} (we are collecting your email to be able to \textbf{send you the links when each Unit is published and give you information about the challenges and updates).}

The best way to keep in touch is to join our discord server to exchange with the community and with us \url{https://discord.gg/ydHrjt3WP5}

\subsection{Prerequisites}\label{prerequisites}

Before diving into the notebook, you need to:

\textbf{Study \href{https://huggingface.co/deep-rl-course/unit2/introduction}{Q-Learning by reading Unit 2}}

\subsection{A small recap of Q-Learning}\label{a-small-recap-of-q-learning}

\emph{Q-Learning} \textbf{is the RL algorithm that}:

\begin{itemize}
\item
  Trains \emph{Q-Function}, an \textbf{action-value function} that is encoded, in internal memory, by a \emph{Q-table} \textbf{that contains all the state-action pair values.}
\item
  Given a state and action, our Q-Function \textbf{will search the Q-table for the corresponding value.}
\end{itemize}

\CourseFigure{Q-function-2.jpg}{Q-table lookup: obtaining the value of a state-action pair \protect\href{https://huggingface.co/datasets/huggingface-deep-rl-course/course-images/resolve/main/en/unit3/Q-function-2.jpg}{[original figure]}}

\begin{itemize}
\item
  When the training is done, \textbf{we have an optimal Q-Function, so an optimal Q-Table.}
\item
  And if we \textbf{have an optimal Q-function}, we
\end{itemize}

have an optimal policy, since we \textbf{know for each state, the best action to take.}

\CourseFigure{link-value-policy.jpg}{Recovering an optimal policy from the optimal action-value function \protect\href{https://huggingface.co/datasets/huggingface-deep-rl-course/course-images/resolve/main/en/unit3/link-value-policy.jpg}{[original figure]}}

But, in the beginning, our \textbf{Q-Table is useless since it gives arbitrary value for each state-action pair (most of the time we initialize the Q-Table to 0 values)}. But, as we'll explore the environment and update our Q-Table it will give us better and better approximations

\CourseFigure{q-learning.jpeg}{Learning a Q-table from zero initialization toward a useful policy \protect\href{https://huggingface.co/datasets/huggingface-deep-rl-course/course-images/resolve/main/en/notebooks/unit2/q-learning.jpeg}{[original figure]}}

This is the Q-Learning pseudocode:

\CourseFigure{Q-learning-2.jpg}{Q-learning: reference algorithm and update pseudocode \protect\href{https://huggingface.co/datasets/huggingface-deep-rl-course/course-images/resolve/main/en/unit3/Q-learning-2.jpg}{[original figure]}}

\section{Let's code our first Reinforcement Learning algorithm}\label{lets-code-our-first-reinforcement-learning-algorithm}

To validate this hands-on for the \href{https://huggingface.co/deep-rl-course/en/unit0/introduction\#certification-process}{certification process}, you need to push your trained Taxi model to the Hub and \textbf{get a result of \textgreater= 4.5}.

To find your result, go to the \href{https://huggingface.co/spaces/huggingface-projects/Deep-Reinforcement-Learning-Leaderboard}{leaderboard} and find your model, \textbf{the result = mean\_reward - std of reward}

For more information about the certification process, check this section \url{https://huggingface.co/deep-rl-course/en/unit0/introduction\#certification-process}

\subsection{Install dependencies and create a virtual display}\label{install-dependencies-and-create-a-virtual-display}

In the notebook, we'll need to generate a replay video. To do so, with Colab, \textbf{we need to have a virtual screen to render the environment} (and thus record the frames).

Hence the following cell will install the libraries and create and run a virtual screen

We'll install multiple ones:

\begin{itemize}
\tightlist
\item
  \passthrough{\lstinline!gymnasium!}: Contains the FrozenLake-v1 and Taxi-v3 environments.
\item
  \passthrough{\lstinline!pygame!}: Used for the FrozenLake-v1 and Taxi-v3 UI.
\item
  \passthrough{\lstinline!numpy!}: Used for handling our Q-table.
\end{itemize}

The Hugging Face Hub works as a central place where anyone can share and explore models and datasets. It has versioning, metrics, visualizations and other features that will allow you to easily collaborate with others.

You can see here all the Deep RL models available (if they use Q Learning) here \url{https://huggingface.co/models?other=q-learning}

\begin{lstlisting}[language=bash]
pip install -r https://raw.githubusercontent.com/huggingface/deep-rl-class/main/notebooks/unit2/requirements-unit2.txt
\end{lstlisting}

\begin{lstlisting}[language=bash]
sudo apt-get update
sudo apt-get install -y python3-opengl
apt install ffmpeg xvfb
pip3 install pyvirtualdisplay
\end{lstlisting}

To make sure the new installed libraries are used, \textbf{sometimes it's required to restart the notebook runtime}. The next cell will force the \textbf{runtime to crash, so you'll need to connect again and run the code starting from here}. Thanks to this trick, \textbf{we will be able to run our virtual screen.}

\begin{lstlisting}[language=Python]
import os

os.kill(os.getpid(), 9)
\end{lstlisting}

\begin{lstlisting}[language=Python]
# Virtual display
from pyvirtualdisplay import Display

virtual_display = Display(visible=0, size=(1400, 900))
virtual_display.start()
\end{lstlisting}

\subsection{Import the packages}\label{import-the-packages}

In addition to the installed libraries, we also use:

\begin{itemize}
\tightlist
\item
  \passthrough{\lstinline!random!}: To generate random numbers (that will be useful for epsilon-greedy policy).
\item
  \passthrough{\lstinline!imageio!}: To generate a replay video.
\end{itemize}

\begin{lstlisting}[language=Python]
import numpy as np
import gymnasium as gym
import random
import imageio
import os
import tqdm

import pickle5 as pickle
from tqdm.notebook import tqdm
\end{lstlisting}

We're now ready to code our Q-Learning algorithm

\section{Part 1: Frozen Lake (non slippery version)}\label{part-1-frozen-lake-non-slippery-version}

\subsection{\texorpdfstring{\#\# Create and understand \href{https://gymnasium.farama.org/environments/toy_text/frozen_lake/}{FrozenLake environment}}{\#\# Create and understand FrozenLake environment}}\label{create-and-understand-frozenlake-environment}

A good habit when you start to use an environment is to check its documentation

\url{https://gymnasium.farama.org/environments/toy_text/frozen_lake/}

\begin{center}\rule{0.5\linewidth}{0.5pt}\end{center}

We're going to train our Q-Learning agent \textbf{to navigate from the starting state (S) to the goal state (G) by walking only on frozen tiles (F) and avoid holes (H)}.

We can have two sizes of environment:

\begin{itemize}
\tightlist
\item
  \passthrough{\lstinline!map_name="4x4"!}: a 4x4 grid version
\item
  \passthrough{\lstinline!map_name="8x8"!}: a 8x8 grid version
\end{itemize}

The environment has two modes:

\begin{itemize}
\tightlist
\item
  \passthrough{\lstinline!is_slippery=False!}: The agent always moves \textbf{in the intended direction} due to the non-slippery nature of the frozen lake (deterministic).
\item
  \passthrough{\lstinline!is_slippery=True!}: The agent \textbf{may not always move in the intended direction} due to the slippery nature of the frozen lake (stochastic).
\end{itemize}

For now let's keep it simple with the 4x4 map and non-slippery. We add a parameter called \passthrough{\lstinline!render_mode!} that specifies how the environment should be visualised. In our case because we \textbf{want to record a video of the environment at the end, we need to set render\_mode to rgb\_array}.

As \href{https://gymnasium.farama.org/api/env/\#gymnasium.Env.render}{explained in the documentation} ``rgb\_array'': Return a single frame representing the current state of the environment. A frame is a np.ndarray with shape (x, y, 3) representing RGB values for an x-by-y pixel image.

\begin{lstlisting}[language=Python]
# Create the FrozenLake-v1 environment using 4x4 map and non-slippery version and render_mode="rgb_array"
env = gym.make()  # TODO use the correct parameters
\end{lstlisting}

\subsubsection{Solution}\label{solution}

\begin{lstlisting}[language=Python]
env = gym.make("FrozenLake-v1", map_name="4x4", is_slippery=False, render_mode="rgb_array")
\end{lstlisting}

You can create your own custom grid like this:

\begin{lstlisting}[language=Python]
desc=["SFFF", "FHFH", "FFFH", "HFFG"]
gym.make('FrozenLake-v1', desc=desc, is_slippery=True)
\end{lstlisting}

but we'll use the default environment for now.

\subsubsection{Let's see what the Environment looks like:}\label{lets-see-what-the-environment-looks-like}

\begin{lstlisting}[language=Python]
# We create our environment with gym.make("<name_of_the_environment>")- `is_slippery=False`: The agent always moves in the intended direction due to the non-slippery nature of the frozen lake (deterministic).
print("_____OBSERVATION SPACE_____ \n")
print("Observation Space", env.observation_space)
print("Sample observation", env.observation_space.sample())  # Get a random observation
\end{lstlisting}

We see with \passthrough{\lstinline!Observation Space Shape Discrete(16)!} that the observation is an integer representing the \textbf{agent's current position as current\_row * ncols + current\_col (where both the row and col start at 0)}.

For example, the goal position in the 4x4 map can be calculated as follows: 3 * 4 + 3 = 15. The number of possible observations is dependent on the size of the map. \textbf{For example, the 4x4 map has 16 possible observations.}

For instance, this is what state = 0 looks like:

\CourseFigure{frozenlake.png}{FrozenLake grid, actions, and tile types \protect\href{https://huggingface.co/datasets/huggingface-deep-rl-course/course-images/resolve/main/en/notebooks/unit2/frozenlake.png}{[original figure]}}

\begin{lstlisting}[language=Python]
print("\n _____ACTION SPACE_____ \n")
print("Action Space Shape", env.action_space.n)
print("Action Space Sample", env.action_space.sample())  # Take a random action
\end{lstlisting}

The action space (the set of possible actions the agent can take) is discrete with 4 actions available :

\begin{itemize}
\tightlist
\item
  0: GO LEFT
\item
  1: GO DOWN
\item
  2: GO RIGHT
\item
  3: GO UP
\end{itemize}

Reward function :

\begin{itemize}
\tightlist
\item
  Reach goal: +1
\item
  Reach hole: 0
\item
  Reach frozen: 0
\end{itemize}

\subsection{Create and Initialize the Q-table}\label{create-and-initialize-the-q-table}

( Step 1 of the pseudocode)

\CourseFigure{Q-learning-2.jpg}{Q-learning: reference algorithm and update pseudocode \protect\href{https://huggingface.co/datasets/huggingface-deep-rl-course/course-images/resolve/main/en/unit3/Q-learning-2.jpg}{[original figure]}}

It's time to initialize our Q-table! To know how many rows (states) and columns (actions) to use, we need to know the action and observation space. We already know their values from before, but we'll want to obtain them programmatically so that our algorithm generalizes for different environments. Gym provides us a way to do that: \passthrough{\lstinline!env.action_space.n!} and \passthrough{\lstinline!env.observation_space.n!}

\begin{lstlisting}[language=Python]
state_space =
print("There are ", state_space, " possible states")

action_space =
print("There are ", action_space, " possible actions")
\end{lstlisting}

\begin{lstlisting}[language=Python]
# Let's create our Qtable of size (state_space, action_space) and initialized each values at 0 using np.zeros. np.zeros needs a tuple (a,b)
def initialize_q_table(state_space, action_space):
  Qtable =
  return Qtable
\end{lstlisting}

\begin{lstlisting}[language=Python]
Qtable_frozenlake = initialize_q_table(state_space, action_space)
\end{lstlisting}

\subsubsection{Solution}\label{solution-1}

\begin{lstlisting}[language=Python]
state_space = env.observation_space.n
print("There are ", state_space, " possible states")

action_space = env.action_space.n
print("There are ", action_space, " possible actions")
\end{lstlisting}

\begin{lstlisting}[language=Python]
# Let's create our Qtable of size (state_space, action_space) and initialized each values at 0 using np.zeros
def initialize_q_table(state_space, action_space):
    Qtable = np.zeros((state_space, action_space))
    return Qtable
\end{lstlisting}

\begin{lstlisting}[language=Python]
Qtable_frozenlake = initialize_q_table(state_space, action_space)
\end{lstlisting}

\subsection{Define the greedy policy}\label{define-the-greedy-policy}

Remember we have two policies since Q-Learning is an \textbf{off-policy} algorithm. This means we're using a \textbf{different policy for acting and updating the value function}.

\begin{itemize}
\tightlist
\item
  Epsilon-greedy policy (acting policy)
\item
  Greedy-policy (updating policy)
\end{itemize}

The greedy policy will also be the final policy we'll have when the Q-learning agent completes training. The greedy policy is used to select an action using the Q-table.

\CourseFigure{off-on-4.jpg}{Comparison of off-policy and on-policy methods \protect\href{https://huggingface.co/datasets/huggingface-deep-rl-course/course-images/resolve/main/en/unit3/off-on-4.jpg}{[original figure]}}

\begin{lstlisting}[language=Python]
def greedy_policy(Qtable, state):
  # Exploitation: take the action with the highest state, action value
  action =

  return action
\end{lstlisting}

\paragraph{Solution}\label{solution-2}

\begin{lstlisting}[language=Python]
def greedy_policy(Qtable, state):
    # Exploitation: take the action with the highest state, action value
    action = np.argmax(Qtable[state][:])

    return action
\end{lstlisting}

\subsection{Define the epsilon-greedy policy}\label{define-the-epsilon-greedy-policy}

Epsilon-greedy is the training policy that handles the exploration/exploitation trade-off.

The idea with epsilon-greedy:

\begin{itemize}
\item
  With \emph{probability 1 - \(\varepsilon\)} : \textbf{we do exploitation} (i.e.~our agent selects the action with the highest state-action pair value).
\item
  With \emph{probability \(\varepsilon\)}: we do \textbf{exploration} (trying a random action).
\end{itemize}

As the training continues, we progressively \textbf{reduce the epsilon value since we will need less and less exploration and more exploitation.}

\CourseFigure{Q-learning-4.jpg}{Epsilon-greedy strategy: exploration and exploitation \protect\href{https://huggingface.co/datasets/huggingface-deep-rl-course/course-images/resolve/main/en/unit3/Q-learning-4.jpg}{[original figure]}}

\begin{lstlisting}[language=Python]
def epsilon_greedy_policy(Qtable, state, epsilon):
  # Randomly generate a number between 0 and 1
  random_num =
  # if random_num > greater than epsilon --> exploitation
  if random_num > epsilon:
    # Take the action with the highest value given a state
    # np.argmax can be useful here
    action =
  # else --> exploration
  else:
    action = # Take a random action

  return action
\end{lstlisting}

\paragraph{Solution}\label{solution-3}

\begin{lstlisting}[language=Python]
def epsilon_greedy_policy(Qtable, state, epsilon):
    # Randomly generate a number between 0 and 1
    random_num = random.uniform(0, 1)
    # if random_num > greater than epsilon --> exploitation
    if random_num > epsilon:
        # Take the action with the highest value given a state
        # np.argmax can be useful here
        action = greedy_policy(Qtable, state)
    # else --> exploration
    else:
        action = env.action_space.sample()

    return action
\end{lstlisting}

\subsection{Define the hyperparameters}\label{define-the-hyperparameters}

The exploration related hyperparamters are some of the most important ones.

\begin{itemize}
\tightlist
\item
  We need to make sure that our agent \textbf{explores enough of the state space} to learn a good value approximation. To do that, we need to have progressive decay of the epsilon.
\item
  If you decrease epsilon too fast (too high decay\_rate), \textbf{you take the risk that your agent will be stuck}, since your agent didn't explore enough of the state space and hence can't solve the problem.
\end{itemize}

\begin{lstlisting}[language=Python]
# Training parameters
n_training_episodes = 10000  # Total training episodes
learning_rate = 0.7  # Learning rate

# Evaluation parameters
n_eval_episodes = 100  # Total number of test episodes

# Environment parameters
env_id = "FrozenLake-v1"  # Name of the environment
max_steps = 99  # Max steps per episode
gamma = 0.95  # Discounting rate
eval_seed = []  # The evaluation seed of the environment

# Exploration parameters
max_epsilon = 1.0  # Exploration probability at start
min_epsilon = 0.05  # Minimum exploration probability
decay_rate = 0.0005  # Exponential decay rate for exploration prob
\end{lstlisting}

\subsection{Create the training loop method}\label{create-the-training-loop-method}

\CourseFigure{Q-learning-2.jpg}{Q-learning: reference algorithm and update pseudocode \protect\href{https://huggingface.co/datasets/huggingface-deep-rl-course/course-images/resolve/main/en/unit3/Q-learning-2.jpg}{[original figure]}}

The training loop goes like this:

\begin{lstlisting}
For episode in the total of training episodes:

Reduce epsilon (since we need less and less exploration)
Reset the environment

  For step in max timesteps:
    Choose the action At using epsilon greedy policy
    Take the action (a) and observe the outcome state(s') and reward (r)
    Update the Q-value Q(s,a) using Bellman equation Q(s,a) + lr [R(s,a) + gamma * max Q(s',a') - Q(s,a)]
    If done, finish the episode
    Our next state is the new state
\end{lstlisting}

\begin{lstlisting}[language=Python]
def train(n_training_episodes, min_epsilon, max_epsilon, decay_rate, env, max_steps, Qtable):
  for episode in tqdm(range(n_training_episodes)):
    # Reduce epsilon (because we need less and less exploration)
    epsilon = min_epsilon + (max_epsilon - min_epsilon)*np.exp(-decay_rate*episode)
    # Reset the environment
    state, info = env.reset()
    step = 0
    terminated = False
    truncated = False

    # repeat
    for step in range(max_steps):
      # Choose the action At using epsilon greedy policy
      action =

      # Take action At and observe Rt+1 and St+1
      # Take the action (a) and observe the outcome state(s') and reward (r)
      new_state, reward, terminated, truncated, info =

      # Update Q(s,a):= Q(s,a) + lr [R(s,a) + gamma * max Q(s',a') - Q(s,a)]
      Qtable[state][action] =

      # If terminated or truncated finish the episode
      if terminated or truncated:
        break

      # Our next state is the new state
      state = new_state
  return Qtable
\end{lstlisting}

\paragraph{Solution}\label{solution-4}

\begin{lstlisting}[language=Python]
def train(n_training_episodes, min_epsilon, max_epsilon, decay_rate, env, max_steps, Qtable):
    for episode in tqdm(range(n_training_episodes)):
        # Reduce epsilon (because we need less and less exploration)
        epsilon = min_epsilon + (max_epsilon - min_epsilon) * np.exp(-decay_rate * episode)
        # Reset the environment
        state, info = env.reset()
        step = 0
        terminated = False
        truncated = False

        # repeat
        for step in range(max_steps):
            # Choose the action At using epsilon greedy policy
            action = epsilon_greedy_policy(Qtable, state, epsilon)

            # Take action At and observe Rt+1 and St+1
            # Take the action (a) and observe the outcome state(s') and reward (r)
            new_state, reward, terminated, truncated, info = env.step(action)

            # Update Q(s,a):= Q(s,a) + lr [R(s,a) + gamma * max Q(s',a') - Q(s,a)]
            Qtable[state][action] = Qtable[state][action] + learning_rate * (
                reward + gamma * np.max(Qtable[new_state]) - Qtable[state][action]
            )

            # If terminated or truncated finish the episode
            if terminated or truncated:
                break

            # Our next state is the new state
            state = new_state
    return Qtable
\end{lstlisting}

\subsection{Train the Q-Learning agent}\label{train-the-q-learning-agent}

\begin{lstlisting}[language=Python]
Qtable_frozenlake = train(n_training_episodes, min_epsilon, max_epsilon, decay_rate, env, max_steps, Qtable_frozenlake)
\end{lstlisting}

\subsection{Let's see what our Q-Learning table looks like now}\label{lets-see-what-our-q-learning-table-looks-like-now}

\begin{lstlisting}[language=Python]
Qtable_frozenlake
\end{lstlisting}

\subsection{The evaluation method}\label{the-evaluation-method}

\begin{itemize}
\tightlist
\item
  We defined the evaluation method that we're going to use to test our Q-Learning agent.
\end{itemize}

\begin{lstlisting}[language=Python]
def evaluate_agent(env, max_steps, n_eval_episodes, Q, seed):
    """
    Evaluate the agent for ``n_eval_episodes`` episodes and returns average reward and std of reward.
    :param env: The evaluation environment
    :param n_eval_episodes: Number of episode to evaluate the agent
    :param Q: The Q-table
    :param seed: The evaluation seed array (for taxi-v3)
    """
    episode_rewards = []
    for episode in tqdm(range(n_eval_episodes)):
        if seed:
            state, info = env.reset(seed=seed[episode])
        else:
            state, info = env.reset()
        step = 0
        truncated = False
        terminated = False
        total_rewards_ep = 0

        for step in range(max_steps):
            # Take the action (index) that have the maximum expected future reward given that state
            action = greedy_policy(Q, state)
            new_state, reward, terminated, truncated, info = env.step(action)
            total_rewards_ep += reward

            if terminated or truncated:
                break
            state = new_state
        episode_rewards.append(total_rewards_ep)
    mean_reward = np.mean(episode_rewards)
    std_reward = np.std(episode_rewards)

    return mean_reward, std_reward
\end{lstlisting}

\subsection{Evaluate our Q-Learning agent}\label{evaluate-our-q-learning-agent}

\begin{itemize}
\tightlist
\item
  Usually, you should have a mean reward of 1.0
\item
  The \textbf{environment is relatively easy} since the state space is really small (16). What you can try to do is \href{https://www.gymlibrary.dev/environments/toy_text/frozen_lake/}{to replace it with the slippery version}, which introduces stochasticity, making the environment more complex.
\end{itemize}

\begin{lstlisting}[language=Python]
# Evaluate our Agent
mean_reward, std_reward = evaluate_agent(env, max_steps, n_eval_episodes, Qtable_frozenlake, eval_seed)
print(f"Mean_reward={mean_reward:.2f} +/- {std_reward:.2f}")
\end{lstlisting}

\subsection{Publish our trained model to the Hub}\label{publish-our-trained-model-to-the-hub}

Now that we saw good results after the training, \textbf{we can publish our trained model to the Hub with one line of code}.

Here's an example of a Model Card:

\CourseFigure{modelcard.png}{Model card \protect\href{https://huggingface.co/datasets/huggingface-deep-rl-course/course-images/resolve/main/en/notebooks/unit2/modelcard.png}{[original figure]}}

Under the hood, the Hub uses git-based repositories (don't worry if you don't know what git is), which means you can update the model with new versions as you experiment and improve your agent.

\paragraph{Do not modify this code}\label{do-not-modify-this-code}

\begin{lstlisting}[language=Python]
from huggingface_hub import HfApi, snapshot_download
from huggingface_hub.repocard import metadata_eval_result, metadata_save

from pathlib import Path
import datetime
import json
\end{lstlisting}

\begin{lstlisting}[language=Python]
def record_video(env, Qtable, out_directory, fps=1):
    """
    Generate a replay video of the agent
    :param env
    :param Qtable: Qtable of our agent
    :param out_directory
    :param fps: how many frame per seconds (with taxi-v3 and frozenlake-v1 we use 1)
    """
    images = []
    terminated = False
    truncated = False
    state, info = env.reset(seed=random.randint(0, 500))
    img = env.render()
    images.append(img)
    while not terminated or truncated:
        # Take the action (index) that have the maximum expected future reward given that state
        action = np.argmax(Qtable[state][:])
        state, reward, terminated, truncated, info = env.step(
            action
        )  # We directly put next_state = state for recording logic
        img = env.render()
        images.append(img)
    imageio.mimsave(out_directory, [np.array(img) for i, img in enumerate(images)], fps=fps)
\end{lstlisting}

\begin{lstlisting}[language=Python]
def push_to_hub(repo_id, model, env, video_fps=1, local_repo_path="hub"):
    """
    Evaluate, Generate a video and Upload a model to Hugging Face Hub.
    This method does the complete pipeline:
    - It evaluates the model
    - It generates the model card
    - It generates a replay video of the agent
    - It pushes everything to the Hub

    :param repo_id: repo_id: id of the model repository from the Hugging Face Hub
    :param env
    :param video_fps: how many frame per seconds to record our video replay
    (with taxi-v3 and frozenlake-v1 we use 1)
    :param local_repo_path: where the local repository is
    """
    _, repo_name = repo_id.split("/")

    eval_env = env
    api = HfApi()

    # Step 1: Create the repo
    repo_url = api.create_repo(
        repo_id=repo_id,
        exist_ok=True,
    )

    # Step 2: Download files
    repo_local_path = Path(snapshot_download(repo_id=repo_id))

    # Step 3: Save the model
    if env.spec.kwargs.get("map_name"):
        model["map_name"] = env.spec.kwargs.get("map_name")
        if env.spec.kwargs.get("is_slippery", "") == False:
            model["slippery"] = False

    # Pickle the model
    with open((repo_local_path) / "q-learning.pkl", "wb") as f:
        pickle.dump(model, f)

    # Step 4: Evaluate the model and build JSON with evaluation metrics
    mean_reward, std_reward = evaluate_agent(
        eval_env, model["max_steps"], model["n_eval_episodes"], model["qtable"], model["eval_seed"]
    )

    evaluate_data = {
        "env_id": model["env_id"],
        "mean_reward": mean_reward,
        "n_eval_episodes": model["n_eval_episodes"],
        "eval_datetime": datetime.datetime.now().isoformat(),
    }

    # Write a JSON file called "results.json" that will contain the
    # evaluation results
    with open(repo_local_path / "results.json", "w") as outfile:
        json.dump(evaluate_data, outfile)

    # Step 5: Create the model card
    env_name = model["env_id"]
    if env.spec.kwargs.get("map_name"):
        env_name += "-" + env.spec.kwargs.get("map_name")

    if env.spec.kwargs.get("is_slippery", "") == False:
        env_name += "-" + "no_slippery"

    metadata = {}
    metadata["tags"] = [env_name, "q-learning", "reinforcement-learning", "custom-implementation"]

    # Add metrics
    eval = metadata_eval_result(
        model_pretty_name=repo_name,
        task_pretty_name="reinforcement-learning",
        task_id="reinforcement-learning",
        metrics_pretty_name="mean_reward",
        metrics_id="mean_reward",
        metrics_value=f"{mean_reward:.2f} +/- {std_reward:.2f}",
        dataset_pretty_name=env_name,
        dataset_id=env_name,
    )

    # Merges both dictionaries
    metadata = {**metadata, **eval}

    model_card = f"""
  # **Q-Learning** Agent playing1 **{env_id}**
  This is a trained model of a **Q-Learning** agent playing **{env_id}** .

  ## Usage

  model = load_from_hub(repo_id="{repo_id}", filename="q-learning.pkl")

  # Don't forget to check if you need to add additional attributes (is_slippery=False etc)
  env = gym.make(model["env_id"])
  """

    evaluate_agent(env, model["max_steps"], model["n_eval_episodes"], model["qtable"], model["eval_seed"])

    readme_path = repo_local_path / "README.md"
    readme = ""
    print(readme_path.exists())
    if readme_path.exists():
        with readme_path.open("r", encoding="utf8") as f:
            readme = f.read()
    else:
        readme = model_card

    with readme_path.open("w", encoding="utf-8") as f:
        f.write(readme)

    # Save our metrics to Readme metadata
    metadata_save(readme_path, metadata)

    # Step 6: Record a video
    video_path = repo_local_path / "replay.mp4"
    record_video(env, model["qtable"], video_path, video_fps)

    # Step 7. Push everything to the Hub
    api.upload_folder(
        repo_id=repo_id,
        folder_path=repo_local_path,
        path_in_repo=".",
    )

    print("Your model is pushed to the Hub. You can view your model here: ", repo_url)
\end{lstlisting}

\subsubsection{.}\label{section}

By using \passthrough{\lstinline!push_to_hub!} \textbf{you evaluate, record a replay, generate a model card of your agent and push it to the Hub}.

This way:

\begin{itemize}
\tightlist
\item
  You can \textbf{showcase our work}
\item
  You can \textbf{visualize your agent playing}
\item
  You can \textbf{share an agent with the community that others can use}
\item
  You can \textbf{access a leaderboard to see how well your agent is performing compared to your classmates} \url{https://huggingface.co/spaces/huggingface-projects/Deep-Reinforcement-Learning-Leaderboard}
\end{itemize}

To be able to share your model with the community there are three more steps to follow:

1 (If it's not already done) create an account to HF \url{https://huggingface.co/join}

2 Sign in and then, you need to store your authentication token from the Hugging Face website.

\begin{itemize}
\tightlist
\item
  Create a new token (\url{https://huggingface.co/settings/tokens}) \textbf{with write role}
\end{itemize}

\CourseFigure{create-token.jpg}{Generating a Hugging Face Hub access token \protect\href{https://huggingface.co/datasets/huggingface-deep-rl-course/course-images/resolve/main/en/notebooks/create-token.jpg}{[original figure]}}

\begin{lstlisting}[language=Python]
from huggingface_hub import notebook_login

notebook_login()
\end{lstlisting}

If you don't want to use a Google Colab or a Jupyter Notebook, you need to use this command instead: \passthrough{\lstinline!huggingface-cli login!} (or \passthrough{\lstinline!login!})

3 We're now ready to push our trained agent to the Hub using \passthrough{\lstinline!push_to_hub()!} function

\begin{itemize}
\tightlist
\item
  Let's create \textbf{the model dictionary that contains the hyperparameters and the Q\_table}.
\end{itemize}

\begin{lstlisting}[language=Python]
model = {
    "env_id": env_id,
    "max_steps": max_steps,
    "n_training_episodes": n_training_episodes,
    "n_eval_episodes": n_eval_episodes,
    "eval_seed": eval_seed,
    "learning_rate": learning_rate,
    "gamma": gamma,
    "max_epsilon": max_epsilon,
    "min_epsilon": min_epsilon,
    "decay_rate": decay_rate,
    "qtable": Qtable_frozenlake,
}
\end{lstlisting}

Let's fill the \passthrough{\lstinline!push_to_hub!} function:

\begin{itemize}
\tightlist
\item
  \passthrough{\lstinline!repo_id!}: the name of the Hugging Face Hub Repository that will be created/updated `
\end{itemize}

(repo\_id = \{username\}/\{repo\_name\})\passthrough{\lstinline!A good!}repo\_id\passthrough{\lstinline!is!}\{username\}/q-\{env\_id\}`

\begin{itemize}
\tightlist
\item
  \passthrough{\lstinline!model!}: our model dictionary containing the hyperparameters and the Qtable.
\item
  \passthrough{\lstinline!env!}: the environment.
\item
  \passthrough{\lstinline!commit_message!}: message of the commit
\end{itemize}

\begin{lstlisting}[language=Python]
model
\end{lstlisting}

\begin{lstlisting}[language=Python]
username = ""  # FILL THIS
repo_name = "q-FrozenLake-v1-4x4-noSlippery"
push_to_hub(repo_id=f"{username}/{repo_name}", model=model, env=env)
\end{lstlisting}

Congrats you've just implemented from scratch, trained, and uploaded your first Reinforcement Learning agent. FrozenLake-v1 no\_slippery is very simple environment, let's try a harder one .

\section{Part 2: Taxi-v3}\label{part-2-taxi-v3}

\subsection{\texorpdfstring{\#\# Create and understand \href{https://gymnasium.farama.org/environments/toy_text/taxi/}{Taxi-v3}}{\#\# Create and understand Taxi-v3}}\label{create-and-understand-taxi-v3}

A good habit when you start to use an environment is to check its documentation

\url{https://gymnasium.farama.org/environments/toy_text/taxi/}

\begin{center}\rule{0.5\linewidth}{0.5pt}\end{center}

In \passthrough{\lstinline!Taxi-v3!} , there are four designated locations in the grid world indicated by R(ed), G(reen), Y(ellow), and B(lue).

When the episode starts, \textbf{the taxi starts off at a random square} and the passenger is at a random location. The taxi drives to the passenger's location, \textbf{picks up the passenger}, drives to the passenger's destination (another one of the four specified locations), and then \textbf{drops off the passenger}. Once the passenger is dropped off, the episode ends.

\CourseFigure{taxi.png}{Taxi environment used for Q-learning \protect\href{https://huggingface.co/datasets/huggingface-deep-rl-course/course-images/resolve/main/en/notebooks/unit2/taxi.png}{[original figure]}}

\begin{lstlisting}[language=Python]
env = gym.make("Taxi-v3", render_mode="rgb_array")
\end{lstlisting}

There are \textbf{500 discrete states since there are 25 taxi positions, 5 possible locations of the passenger} (including the case when the passenger is in the taxi), and \textbf{4 destination locations.}

\begin{lstlisting}[language=Python]
state_space = env.observation_space.n
print("There are ", state_space, " possible states")
\end{lstlisting}

\begin{lstlisting}[language=Python]
action_space = env.action_space.n
print("There are ", action_space, " possible actions")
\end{lstlisting}

The action space (the set of possible actions the agent can take) is discrete with \textbf{6 actions available }:

\begin{itemize}
\tightlist
\item
  0: move south
\item
  1: move north
\item
  2: move east
\item
  3: move west
\item
  4: pickup passenger
\item
  5: drop off passenger
\end{itemize}

Reward function :

\begin{itemize}
\tightlist
\item
  -1 per step unless other reward is triggered.
\item
  +20 delivering passenger.
\item
  -10 executing ``pickup'' and ``drop-off'' actions illegally.
\end{itemize}

\begin{lstlisting}[language=Python]
# Create our Q table with state_size rows and action_size columns (500x6)
Qtable_taxi = initialize_q_table(state_space, action_space)
print(Qtable_taxi)
print("Q-table shape: ", Qtable_taxi.shape)
\end{lstlisting}

\subsection{Define the hyperparameters}\label{define-the-hyperparameters-1}

DO NOT MODIFY EVAL\_SEED: the eval\_seed array \textbf{allows us to evaluate your agent with the same taxi starting positions for every classmate}

\begin{lstlisting}[language=Python]
# Training parameters
n_training_episodes = 25000  # Total training episodes
learning_rate = 0.7  # Learning rate

# Evaluation parameters
n_eval_episodes = 100  # Total number of test episodes

# DO NOT MODIFY EVAL_SEED
eval_seed = [
    16,
    54,
    165,
    177,
    191,
    191,
    120,
    80,
    149,
    178,
    48,
    38,
    6,
    125,
    174,
    73,
    50,
    172,
    100,
    148,
    146,
    6,
    25,
    40,
    68,
    148,
    49,
    167,
    9,
    97,
    164,
    176,
    61,
    7,
    54,
    55,
    161,
    131,
    184,
    51,
    170,
    12,
    120,
    113,
    95,
    126,
    51,
    98,
    36,
    135,
    54,
    82,
    45,
    95,
    89,
    59,
    95,
    124,
    9,
    113,
    58,
    85,
    51,
    134,
    121,
    169,
    105,
    21,
    30,
    11,
    50,
    65,
    12,
    43,
    82,
    145,
    152,
    97,
    106,
    55,
    31,
    85,
    38,
    112,
    102,
    168,
    123,
    97,
    21,
    83,
    158,
    26,
    80,
    63,
    5,
    81,
    32,
    11,
    28,
    148,
]  # Evaluation seed, this ensures that all classmates agents are trained on the same taxi starting position
# Each seed has a specific starting state
# Environment parameters
env_id = "Taxi-v3"  # Name of the environment
max_steps = 99  # Max steps per episode
gamma = 0.95  # Discounting rate

# Exploration parameters
max_epsilon = 1.0  # Exploration probability at start
min_epsilon = 0.05  # Minimum exploration probability
decay_rate = 0.005  # Exponential decay rate for exploration prob
\end{lstlisting}

\subsection{Train our Q-Learning agent}\label{train-our-q-learning-agent}

\begin{lstlisting}[language=Python]
Qtable_taxi = train(n_training_episodes, min_epsilon, max_epsilon, decay_rate, env, max_steps, Qtable_taxi)
Qtable_taxi
\end{lstlisting}

\subsection{Create a model dictionary and publish our trained model to the Hub}\label{create-a-model-dictionary-and-publish-our-trained-model-to-the-hub}

\begin{itemize}
\tightlist
\item
  We create a model dictionary that will contain all the training hyperparameters for reproducibility and the Q-Table.
\end{itemize}

\begin{lstlisting}[language=Python]
model = {
    "env_id": env_id,
    "max_steps": max_steps,
    "n_training_episodes": n_training_episodes,
    "n_eval_episodes": n_eval_episodes,
    "eval_seed": eval_seed,
    "learning_rate": learning_rate,
    "gamma": gamma,
    "max_epsilon": max_epsilon,
    "min_epsilon": min_epsilon,
    "decay_rate": decay_rate,
    "qtable": Qtable_taxi,
}
\end{lstlisting}

\begin{lstlisting}[language=Python]
username = ""  # FILL THIS
repo_name = ""  # FILL THIS
push_to_hub(repo_id=f"{username}/{repo_name}", model=model, env=env)
\end{lstlisting}

Now that it's on the Hub, you can compare the results of your Taxi-v3 with your classmates using the leaderboard \url{https://huggingface.co/spaces/huggingface-projects/Deep-Reinforcement-Learning-Leaderboard}

\CourseFigure{taxi-leaderboard.png}{Taxi Leaderboard \protect\href{https://huggingface.co/datasets/huggingface-deep-rl-course/course-images/resolve/main/en/notebooks/unit2/taxi-leaderboard.png}{[original figure]}}

\section{Part 3: Load from Hub}\label{part-3-load-from-hub}

What's amazing with Hugging Face Hub is that you can easily load powerful models from the community.

Loading a saved model from the Hub is really easy:

\begin{enumerate}
\def\labelenumi{\arabic{enumi}.}
\tightlist
\item
  You go \url{https://huggingface.co/models?other=q-learning} to see the list of all the q-learning saved models.
\item
  You select one and copy its repo\_id
\end{enumerate}

\CourseFigure{copy-id.png}{Example of locating a model repository ID on Hugging Face Hub \protect\href{https://huggingface.co/datasets/huggingface-deep-rl-course/course-images/resolve/main/en/notebooks/unit2/copy-id.png}{[original figure]}}

\begin{enumerate}
\def\labelenumi{\arabic{enumi}.}
\setcounter{enumi}{2}
\tightlist
\item
  Then we just need to use \passthrough{\lstinline!load_from_hub!} with:
\end{enumerate}

\begin{itemize}
\tightlist
\item
  The repo\_id
\item
  The filename: the saved model inside the repo.
\end{itemize}

\paragraph{Do not modify this code}\label{do-not-modify-this-code-1}

\begin{lstlisting}[language=Python]
from urllib.error import HTTPError

from huggingface_hub import hf_hub_download


def load_from_hub(repo_id: str, filename: str) -> str:
    """
    Download a model from Hugging Face Hub.
    :param repo_id: id of the model repository from the Hugging Face Hub
    :param filename: name of the model zip file from the repository
    """
    # Get the model from the Hub, download and cache the model on your local disk
    pickle_model = hf_hub_download(repo_id=repo_id, filename=filename)

    with open(pickle_model, "rb") as f:
        downloaded_model_file = pickle.load(f)

    return downloaded_model_file
\end{lstlisting}

\subsubsection{.}\label{section-1}

\begin{lstlisting}[language=Python]
model = load_from_hub(repo_id="ThomasSimonini/q-Taxi-v3", filename="q-learning.pkl")  # Try to use another model

print(model)
env = gym.make(model["env_id"])

evaluate_agent(env, model["max_steps"], model["n_eval_episodes"], model["qtable"], model["eval_seed"])
\end{lstlisting}

\begin{lstlisting}[language=Python]
model = load_from_hub(
    repo_id="ThomasSimonini/q-FrozenLake-v1-no-slippery", filename="q-learning.pkl"
)  # Try to use another model

env = gym.make(model["env_id"], is_slippery=False)

evaluate_agent(env, model["max_steps"], model["n_eval_episodes"], model["qtable"], model["eval_seed"])
\end{lstlisting}

\subsection{Some additional challenges}\label{some-additional-challenges}

The best way to learn \textbf{is to try things on your own}! As you saw, the current agent is not doing great. As a first suggestion, you can train for more steps. With 1,000,000 steps, we saw some great results!

In the \href{https://huggingface.co/spaces/huggingface-projects/Deep-Reinforcement-Learning-Leaderboard}{Leaderboard} you will find your agents. Can you get to the top?

Here are some ideas to climb up the leaderboard:

\begin{itemize}
\tightlist
\item
  Train more steps
\item
  Try different hyperparameters by looking at what your classmates have done.
\item
  \textbf{Push your new trained model} on the Hub
\end{itemize}

Are walking on ice and driving taxis too boring to you? Try to \textbf{change the environment}, why not use FrozenLake-v1 slippery version? Check how they work \href{https://gymnasium.farama.org/}{using the gymnasium documentation} and have fun .

\begin{center}\rule{0.5\linewidth}{0.5pt}\end{center}

Congrats , you've just implemented, trained, and uploaded your first Reinforcement Learning agent.

Understanding Q-Learning is an \textbf{important step to understanding value-based methods.}

In the next Unit with Deep Q-Learning, we'll see that while creating and updating a Q-table was a good strategy --- \textbf{however, it is not scalable.}

For instance, imagine you create an agent that learns to play Doom.

\CourseFigure{basic.png}{Doom \protect\href{https://vizdoom.cs.put.edu.pl/user/pages/01.tutorial/basic.png}{[original figure]}}

Doom is a large environment with a huge state space (millions of different states). Creating and updating a Q-table for that environment would not be efficient.

That's why we'll study Deep Q-Learning in the next unit, an algorithm \textbf{where we use a neural network that approximates, given a state, the different Q-values for each action.}

\CourseFigure{atari-envs.png}{atari envs \protect\href{https://huggingface.co/datasets/huggingface-deep-rl-course/course-images/resolve/main/en/unit4/atari-envs.gif}{[original animation]}}

See you in Unit 3!

\subsection{Keep learning, stay awesome}\label{keep-learning-stay-awesome}

\section{Second Quiz}\label{second-quiz}

The best way to learn and \href{https://www.coursera.org/lecture/learning-how-to-learn/illusions-of-competence-BuFzf}{to avoid the illusion of competence} \textbf{is to test yourself.} This will help you to find \textbf{where you need to reinforce your knowledge}.

\subsubsection{Q1: What is Q-Learning?}\label{q1-what-is-q-learning}

\textbf{Select all applicable answers.}

\begin{itemize}
\tightlist
\item
  \textbf{A.} The algorithm we use to train our Q-function
\item
  \textbf{B.} A value function
\item
  \textbf{C.} An algorithm that determines the value of being at a particular state and taking a specific action at that state
\item
  \textbf{D.} A table
\end{itemize}

\subsubsection{Q2: What is a Q-table?}\label{q2-what-is-a-q-table}

\textbf{Select all applicable answers.}

\begin{itemize}
\tightlist
\item
  \textbf{A.} An algorithm we use in Q-Learning
\item
  \textbf{B.} Q-table is the internal memory of our agent
\item
  \textbf{C.} In Q-table each cell corresponds a state value
\end{itemize}

\subsubsection{\texorpdfstring{Q3: Why if we have an optimal Q-function \(Q^*\) we have an optimal policy?}{Q3: Why if we have an optimal Q-function Q\^{}* we have an optimal policy?}}\label{q3-why-if-we-have-an-optimal-q-function-q-we-have-an-optimal-policy}

\subsubsection{Q4: Can you explain what is Epsilon-Greedy Strategy?}\label{q4-can-you-explain-what-is-epsilon-greedy-strategy}

\subsubsection{Q5: How do we update the Q value of a state, action pair?}\label{q5-how-do-we-update-the-q-value-of-a-state-action-pair}

\CourseFigure{q-update-ex.jpg}{Quiz: identify terms in the Q-learning update rule \protect\href{https://huggingface.co/datasets/huggingface-deep-rl-course/course-images/resolve/main/en/unit3/q-update-ex.jpg}{[original figure]}}

\subsubsection{Q6: What's the difference between on-policy and off-policy}\label{q6-whats-the-difference-between-on-policy-and-off-policy}

\subsection{Answer Key and Explanations}\label{answer-key-and-explanations-1}

\subsubsection{Q1: What is Q-Learning?}\label{q1-what-is-q-learning-1}

\textbf{Correct choice(s): A}

\begin{itemize}
\item
  \textbf{Option B:} It's an action-value function since it determines the value of being at a particular state and taking a specific action at that state
\item
  \textbf{Option C:} Q-function is the function that determines the value of being at a particular state and taking a specific action at that state.
\item
  \textbf{Option D:} Q-learning is not a Q-table. The Q-function is the algorithm that will feed the Q-table.
\end{itemize}

\subsubsection{Q2: What is a Q-table?}\label{q2-what-is-a-q-table-1}

\textbf{Correct choice(s): B}

\begin{itemize}
\tightlist
\item
  \textbf{Option C:} Each cell corresponds to a state-action value pair value. Not a state value.
\end{itemize}

\subsubsection{\texorpdfstring{Q3: Why if we have an optimal Q-function \(Q^*\) we have an optimal policy?}{Q3: Why if we have an optimal Q-function Q\^{}* we have an optimal policy?}}\label{q3-why-if-we-have-an-optimal-q-function-q-we-have-an-optimal-policy-1}

Because if we have an optimal Q-function, we have an optimal policy since we know for each state what is the best action to take.

\CourseFigure{link-value-policy.jpg}{Recovering an optimal policy from the optimal action-value function \protect\href{https://huggingface.co/datasets/huggingface-deep-rl-course/course-images/resolve/main/en/unit3/link-value-policy.jpg}{[original figure]}}

\subsubsection{Q4: Can you explain what is Epsilon-Greedy Strategy?}\label{q4-can-you-explain-what-is-epsilon-greedy-strategy-1}

Epsilon Greedy Strategy is a policy that handles the exploration/exploitation trade-off.

The idea is that we define epsilon \(\varepsilon\) = 1.0:

\begin{itemize}
\tightlist
\item
  With \emph{probability 1 --- \(\varepsilon\)} : we do exploitation (aka our agent selects the action with the highest state-action pair value).
\item
  With \emph{probability \(\varepsilon\)} : we do exploration (trying random action).
\end{itemize}

\CourseFigure{Q-learning-4.jpg}{Epsilon-greedy strategy: exploration and exploitation \protect\href{https://huggingface.co/datasets/huggingface-deep-rl-course/course-images/resolve/main/en/unit3/Q-learning-4.jpg}{[original figure]}}

\subsubsection{Q5: How do we update the Q value of a state, action pair?}\label{q5-how-do-we-update-the-q-value-of-a-state-action-pair-1}

\CourseFigure{q-update-solution.jpg}{Solution: terms in the Q-learning update rule \protect\href{https://huggingface.co/datasets/huggingface-deep-rl-course/course-images/resolve/main/en/unit3/q-update-solution.jpg}{[original figure]}}

\subsubsection{Q6: What's the difference between on-policy and off-policy}\label{q6-whats-the-difference-between-on-policy-and-off-policy-1}

\CourseFigure{off-on-4.jpg}{Comparison of off-policy and on-policy methods \protect\href{https://huggingface.co/datasets/huggingface-deep-rl-course/course-images/resolve/main/en/unit3/off-on-4.jpg}{[original figure]}}

Congrats on finishing this Quiz , if you missed some elements, take time to read again the chapter to reinforce () your knowledge.

\section{Glossary}\label{glossary}

This is a community-created glossary. Contributions are welcomed!

\subsubsection{Strategies to find the optimal policy}\label{strategies-to-find-the-optimal-policy}

\begin{itemize}
\tightlist
\item
  \textbf{Policy-based methods.} The policy is usually trained with a neural network to select what action to take given a state. In this case it is the neural network which outputs the action that the agent should take instead of using a value function. Depending on the experience received by the environment, the neural network will be re-adjusted and will provide better actions.
\item
  \textbf{Value-based methods.} In this case, a value function is trained to output the value of a state or a state-action pair that will represent our policy. However, this value doesn't define what action the agent should take. In contrast, we need to specify the behavior of the agent given the output of the value function. For example, we could decide to adopt a policy to take the action that always leads to the biggest reward (Greedy Policy). In summary, the policy is a Greedy Policy (or whatever decision the user takes) that uses the values of the value-function to decide the actions to take.
\end{itemize}

\subsubsection{Among the value-based methods, we can find two main strategies}\label{among-the-value-based-methods-we-can-find-two-main-strategies}

\begin{itemize}
\tightlist
\item
  \textbf{The state-value function.} For each state, the state-value function is the expected return if the agent starts in that state and follows the policy until the end.
\item
  \textbf{The action-value function.} In contrast to the state-value function, the action-value calculates for each state and action pair the expected return if the agent starts in that state, takes that action, and then follows the policy forever after.
\end{itemize}

\subsubsection{Epsilon-greedy strategy:}\label{epsilon-greedy-strategy}

\begin{itemize}
\tightlist
\item
  Common strategy used in reinforcement learning that involves balancing exploration and exploitation.
\item
  Chooses the action with the highest expected reward with a probability of 1-epsilon.
\item
  Chooses a random action with a probability of epsilon.
\item
  Epsilon is typically decreased over time to shift focus towards exploitation.
\end{itemize}

\subsubsection{Greedy strategy:}\label{greedy-strategy}

\begin{itemize}
\tightlist
\item
  Involves always choosing the action that is expected to lead to the highest reward, based on the current knowledge of the environment. (Only exploitation)
\item
  Always chooses the action with the highest expected reward.
\item
  Does not include any exploration.
\item
  Can be disadvantageous in environments with uncertainty or unknown optimal actions.
\end{itemize}

\subsubsection{Off-policy vs on-policy algorithms}\label{off-policy-vs-on-policy-algorithms}

\begin{itemize}
\tightlist
\item
  \textbf{Off-policy algorithms:} A different policy is used at training time and inference time
\item
  \textbf{On-policy algorithms:} The same policy is used during training and inference
\end{itemize}

\subsubsection{Monte Carlo and Temporal Difference learning strategies}\label{monte-carlo-and-temporal-difference-learning-strategies}

\begin{itemize}
\item
  \textbf{Monte Carlo (MC):} Learning at the end of the episode. With Monte Carlo, we wait until the episode ends and then we update the value function (or policy function) from a complete episode.
\item
  \textbf{Temporal Difference (TD):} Learning at each step. With Temporal Difference Learning, we update the value function (or policy function) at each step without requiring a complete episode.
\end{itemize}

If you want to improve the course, you can \href{https://github.com/huggingface/deep-rl-class/pulls}{open a Pull Request.}

This glossary was made possible thanks to:

\begin{itemize}
\tightlist
\item
  \href{https://github.com/ramon-rd}{Ramón Rueda}
\item
  \href{https://github.com/hasarinduperera/}{Hasarindu Perera}
\item
  \href{https://github.com/arkadyark/}{Arkady Arkhangorodsky}
\end{itemize}

\section{Conclusion}\label{conclusion}

Congrats on finishing this chapter! There was a lot of information. And congrats on finishing the tutorials. You've just implemented your first RL agent from scratch and shared it on the Hub .

Implementing from scratch when you study a new architecture \textbf{is important to understand how it works.}

It's \textbf{normal if you still feel confused} by all these elements. \textbf{This was the same for me and for everyone who studies RL.}

Take time to really grasp the material before continuing.

In the next chapter, we're going to dive deeper by studying our first Deep Reinforcement Learning algorithm based on Q-Learning: Deep Q-Learning. And you'll train a \textbf{DQN agent with \href{https://github.com/DLR-RM/rl-baselines3-zoo}{RL-Baselines3 Zoo} to play Atari Games}.

\CourseFigure{atari-envs.png}{Atari environments \protect\href{https://huggingface.co/datasets/huggingface-deep-rl-course/course-images/resolve/main/en/unit4/atari-envs.gif}{[original animation]}}

Finally, we would love \textbf{to hear what you think of the course and how we can improve it}. If you have some feedback then please \href{https://forms.gle/BzKXWzLAGZESGNaE9}{fill this form}

\subsubsection{Keep Learning, stay awesome}\label{keep-learning-stay-awesome-1}

\section{Additional Readings}\label{additional-readings}

These are \textbf{optional readings} if you want to go deeper.

\subsection{Monte Carlo and TD Learning}\label{monte-carlo-and-td-learning}

To dive deeper into Monte Carlo and Temporal Difference Learning:

\begin{itemize}
\tightlist
\item
  \href{https://stats.stackexchange.com/questions/355820/why-do-temporal-difference-td-methods-have-lower-variance-than-monte-carlo-met}{Why do temporal difference (TD) methods have lower variance than Monte Carlo methods?}
\item
  \href{https://stats.stackexchange.com/questions/336974/when-are-monte-carlo-methods-preferred-over-temporal-difference-ones}{When are Monte Carlo methods preferred over temporal difference ones?}
\end{itemize}

\subsection{Q-Learning}\label{q-learning}

\begin{itemize}
\tightlist
\item
  \href{http://incompleteideas.net/book/RLbook2020.pdf}{Reinforcement Learning: An Introduction, Richard Sutton and Andrew G. Barto Chapter 5, 6 and 7}
\item
  \href{https://youtu.be/Psrhxy88zww}{Foundations of Deep RL Series, L2 Deep Q-Learning by Pieter Abbeel}
\end{itemize}

\end{document}
